\setgebiet{2}{Zahlentheorie}{mittel}
\chapter{Lösungen: Zahlentheorie}

%%%%%%%%%% 1.1
\stupserzwei{tief}{$\frac{n(n+1)(2n+1)}6 + (n+1)^{2} = \frac{(n+1)\bigl[n(2n+1) + 6(n+1)\bigr]}6$, und die eckige Klammer ist $2n^{2} + 7n + 6 = (n + 2)(2n + 3)$.}{22pt}
\begin{loesungkarte}{tief}{1.1}{Die Summe der Quadrate}
\derweg{Anfang prüfen, im Schritt $(n + 1)$ ausklammern und die verbleibende quadratische Klammer faktorisieren.}
\schritt{(a)}
$n = 1$: $1 = \frac{1\cdot2\cdot3}6$. $n = 2$: $5 = \frac{2\cdot3\cdot5}6$. $n = 3$: $14 = \frac{3\cdot4\cdot7}6$. $n = 10$: $385 = \frac{10\cdot11\cdot21}6 = \frac{2310}6$. \checkmark
\schritt{(b)}
Anfang $n = 1$ erledigt. Schritt: Gelte die Formel für $n$. Dann
\begin{align*} \sum_{k=1}^{n+1}k^{2} &= \frac{n(n+1)(2n+1)}6 + (n+1)^{2}\\ &= \frac{(n+1)\bigl[n(2n+1) + 6(n+1)\bigr]}6 = \frac{(n+1)(2n^{2} + 7n + 6)}6\\ &= \frac{(n+1)(n+2)(2n+3)}6, \end{align*}
und das ist die Formel für $n + 1$, denn $2(n+1) + 1 = 2n + 3$. $\blacksquare$
\schritt{(c)}
Quadrate der Seitenlänge $k$ auf einem $8\times8$-Brett: $(9 - k)^{2}$ Positionen. Summe $8^{2} + 7^{2} + \dots + 1^{2} = \frac{8\cdot9\cdot17}6 = 204$.
\probe{(c) direkt: $64 + 49 + 36 + 25 + 16 + 9 + 4 + 1 = 204$. \checkmark}
\begin{mitnehmen}\kastentext
Im Induktionsschritt nie ausmultiplizieren, was man ausklammern kann: Der Faktor $(n + 1)$ steht in beiden Summanden. Danach ist die Faktorisierung $2n^{2} + 7n + 6 = (n + 2)(2n + 3)$ durch Vieta zu finden --- oder man multipliziert die Zielformel aus und vergleicht.
\end{mitnehmen}
\end{loesungkarte}

%%%%%%%%%% 1.2
\stupserzwei{akzent}{Bei $n = 1\to2$ sind die \enquote{ersten $n$} das Pferd $A$ und die \enquote{letzten $n$} das Pferd $B$. Die Überlappung ist leer.}{22pt}
\begin{loesungkarte}{akzent}{1.2}{Alle Pferde sind gleichfarbig}
\derweg{Der Induktionsschritt wird an der kleinsten Stelle durchgespielt, an der er gebraucht wird --- und dort bricht er.}
\schritt{(a)}
Zwei Pferde $\{A, B\}$. Die ersten $n = 1$: $\{A\}$, gleichfarbig. Die letzten $n = 1$: $\{B\}$, gleichfarbig. Die Überlappung $\{A\}\cap\{B\}$ ist leer --- es gibt kein gemeinsames Pferd, das die Farben verbindet. Der Schluss \enquote{also haben $A$ und $B$ dieselbe Farbe} hat keine Grundlage.
\schritt{(b)}
Für $n\ge2$ ist der Schritt korrekt: Unter $n + 1\ge3$ Pferden haben die ersten $n$ und die letzten $n$ mindestens $n - 1\ge1$ gemeinsame Pferde. Wäre also \emph{die Aussage für $n = 2$} wahr, wären tatsächlich alle Pferde gleichfarbig. Sie ist es nicht.
\schritt{(c)}
Der Anfang $n = 1$ ist richtig, der Schritt gilt ab $n = 2$ --- die Kette hat ein Loch zwischen $1$ und $2$. Ein Induktionsbeweis muss den Schritt für \emph{jedes} $n$ ab dem Anfang liefern; die Frage \enquote{ab welchem $n$ gilt mein Argument wirklich?} gehört zu jedem Beweis.
\schritt{(d)}
Der Schritt behauptet $n\le1\Rightarrow n + 1\le1$, und das ist für kein $n$ wahr. Das \enquote{denn \dots} verdeckt, dass es nichts zu begründen gibt. Hier fehlt nicht der Anfang, sondern der Schritt vollständig.
\begin{fehler}\kastentext
Ein Induktionsschritt, der \enquote{offensichtlich für alle $n$} gilt, ist zu prüfen, wo $n$ am kleinsten ist. Dort stehen Mengen mit einem Element, leere Überlappungen, Summen ohne Summanden --- und dort scheitern die meisten falschen Induktionen.
\end{fehler}
\begin{mitnehmen}\kastentext
Induktion ist Anfang \emph{plus} Schritt \emph{für jedes} $n$ ab dem Anfang. Fehlt der Schritt an einer einzigen Stelle, ist die Kette gerissen --- egal wie gut sie danach aussieht.
\end{mitnehmen}
\end{loesungkarte}

%%%%%%%%%% 1.3
\stupserzwei{gruen}{$S = 1 + 2 + \dots + n$ und $S = n + (n - 1) + \dots + 1$ untereinander: jede Spalte ergibt $n + 1$, es sind $n$ Spalten.}{22pt}
\begin{loesungkarte}{gruen}{1.3}{Die Gaußsche Summe}
\derweg{Zweimal die Summe, einmal umgedreht, spaltenweise addiert; dann die Induktion als formaler Beweis.}
\schritt{(a)}
$2S = (1 + n) + (2 + (n - 1)) + \dots + (n + 1) = n(n + 1)$, also $S = \frac{n(n+1)}2$.
\schritt{(b)}
Anfang $n = 1$: $1 = \frac{1\cdot2}2$. Schritt: $\frac{n(n+1)}2 + (n + 1) = (n + 1)\bigl(\frac n2 + 1\bigr) = \frac{(n+1)(n+2)}2$. $\blacksquare$
\schritt{(c)}
Der Umkehrtrick \emph{findet} die Formel: Man muss sie vorher nicht kennen. Die Induktion \emph{beweist} eine Formel, die man schon hat --- sie ist ein Prüfverfahren, kein Entdeckungsverfahren. Beides braucht man: Entdecken ohne Beweis ist Vermutung, Beweis ohne Entdecken ist unmöglich.
\probe{$n = 100$: $5050$ --- Gauß' Schulaufgabe. \checkmark}
\begin{mitnehmen}\kastentext
Induktion beweist, was man vermutet. Woher die Vermutung kommt, ist eine andere Kunst: kleine Fälle, ein Bild, ein Trick. Die Aufgabensammlung fragt deshalb oft zuerst \enquote{prüfen Sie für $n = 1, 2, 3$} und erst dann \enquote{beweisen Sie}.
\end{mitnehmen}
\end{loesungkarte}

%%%%%%%%%% 1.4
\stupserzwei{gruen}{Ist $n$ nicht prim, so $n = ab$ mit $2\le a, b < n$. Für $a$ und $b$ gilt die Behauptung nach Annahme --- also auch für ihr Produkt.}{22pt}
\begin{loesungkarte}{gruen}{1.4}{Jede Zahl ist ein Produkt von Primzahlen}
\derweg{Starke Induktion: Die Behauptung für alle kleineren Zahlen wird vorausgesetzt, nicht nur für die Vorgängerin.}
\schritt{(a)}
Anfang $n = 2$: prim, Produkt aus einem Faktor. Schritt: Sei $n\ge3$, und jede Zahl $m$ mit $2\le m < n$ sei Produkt von Primzahlen. Ist $n$ prim, fertig. Sonst hat $n$ einen Teiler $a$ mit $1 < a < n$, also $n = ab$ mit $1 < a, b < n$. Nach Voraussetzung sind $a$ und $b$ Produkte von Primzahlen; also ist $n = ab$ eines. $\blacksquare$
\schritt{(b)}
Die Faktoren $a, b$ von $n$ sind irgendwelche kleineren Zahlen --- nicht $n - 1$. Von $n - 1$ allein weiß man nichts über $n$ ($n - 1$ und $n$ sind teilerfremd). Man braucht die Voraussetzung für \emph{alle} Zahlen unter $n$: Das ist die starke Induktion, und sie ist der gewöhnlichen äquivalent (Band 1, Kapitel 1), nur bequemer formuliert.
\schritt{(c)}
Bewiesen ist die \emph{Existenz} einer Primfaktorzerlegung. Es fehlt die \emph{Eindeutigkeit}: dass zwei Zerlegungen bis auf die Reihenfolge übereinstimmen. Sie braucht das Lemma von Euklid ($p\mid ab\Rightarrow p\mid a$ oder $p\mid b$), das aus Bézout folgt (Band 1, Kapitel 3--4). Existenz ist Induktion, Eindeutigkeit ist der ggT.
\begin{mitnehmen}\kastentext
Starke Induktion, wann immer der Schritt auf \enquote{irgendwelche kleineren} zurückgreift: Zerlegungen, Rekursionen, Spiele, Algorithmen. Und: Existenz und Eindeutigkeit sind zwei Sätze mit zwei Beweisen.
\end{mitnehmen}
\end{loesungkarte}

%%%%%%%%%% 1.5
\stupserzwei{gruen}{Drei aufeinanderfolgende Zahlen enthalten ein Vielfaches von $3$ und mindestens ein Vielfaches von $2$; $2$ und $3$ sind teilerfremd, also teilt $6$.}{22pt}
\begin{loesungkarte}{gruen}{1.5}{$n^{3} - n$}
\derweg{Faktorisieren, Reste betrachten, teilerfremde Teiler zusammensetzen.}
\schritt{(a)}
$n^{3} - n = n(n^{2} - 1) = (n - 1)n(n + 1)$. Von drei aufeinanderfolgenden ganzen Zahlen ist genau eine durch $3$ teilbar und mindestens eine gerade. Also $2\mid$ und $3\mid$ Produkt; da $\operatorname{ggT}(2, 3) = 1$, folgt $6\mid n^{3} - n$ (Aufgabe 1.9c).
\schritt{(b)}
Nein: $n = 2$ gibt $8 - 2 = 6$, nicht durch $12$ teilbar. Für $4\mid n^{3} - n$ bräuchte man zwei Faktoren $2$, und unter drei aufeinanderfolgenden Zahlen kann es nur eine gerade geben ($n$ gerade mit $n\equiv2\pmod4$).
\schritt{(c)}
$n^{5} - n = n(n^{4} - 1) = n(n^{2} - 1)(n^{2} + 1) = (n - 1)n(n + 1)(n^{2} + 1)$. Durch $6$ wie in (a). Durch $5$: Ist $n\equiv0, \pm1\pmod5$, teilt $5$ einen der ersten drei Faktoren; ist $n\equiv\pm2$, so $n^{2} + 1\equiv4 + 1\equiv0$. Also $5\mid n^{5} - n$, und mit $\operatorname{ggT}(6, 5) = 1$: $30\mid n^{5} - n$. (Kapitel 7: Fermat gibt $n^{5}\equiv n\pmod5$ in einer Zeile.)
\probe{$n = 3$: $27 - 3 = 24 = 4\cdot6$. \checkmark\ $n = 3$ in (c): $243 - 3 = 240 = 8\cdot30$. \checkmark}
\begin{mitnehmen}\kastentext
Teilbarkeit eines Polynomwerts: faktorisieren, dann Reste durchgehen. Dass $n^{p} - n$ durch $p$ teilbar ist, ist der kleine Fermat --- hier für $p = 2, 3, 5$ mit den Händen.
\end{mitnehmen}
\end{loesungkarte}

%%%%%%%%%% 1.6
\stupserzwei{gruen}{Das Vielfache von $b$ muss \emph{unter} $a$ liegen: für $a = -100$ ist das $-105 = -15\cdot7$, Rest $5$.}{22pt}
\begin{loesungkarte}{gruen}{1.6}{Division mit Rest, dreimal}
\derweg{$q$ ist die größte ganze Zahl mit $qb\le a$; dann $r = a - qb$.}
\schritt{(a), (b), (c)}
$100 = 14\cdot7 + 2$. \quad $-100 = (-15)\cdot7 + 5$ (nicht $(-14)\cdot7 - 2$: der Rest muss in $\{0, \dots, 6\}$ liegen). \quad $7 = 0\cdot100 + 7$.
\schritt{(d)}
Viele Programmiersprachen (C, Java, JavaScript) runden den Quotienten gegen null: $-100 / 7 = -14$, Rest $-100 - (-98) = -2$. Python und die Mathematik runden ab: $-15$, Rest $5$. Beide erfüllen $a = qb + r$; nur die zweite erfüllt $0\le r < b$. Wer Reste negativer Zahlen programmiert, prüft die Konvention --- die Kongruenz $-100\equiv5\equiv-2\pmod7$ ist in beiden Fällen dieselbe.
\probe{$(-15)\cdot7 + 5 = -105 + 5 = -100$. \checkmark}
\begin{mitnehmen}\kastentext
Der Rest ist nie negativ: Das Vielfache liegt immer links von $a$ auf der Zahlengeraden. Für die Zahlentheorie zählt ohnehin nur die Restklasse; $5$ und $-2$ sind modulo $7$ dasselbe.
\end{mitnehmen}
\end{loesungkarte}

%%%%%%%%%% 1.7
\stupserzwei{gruen}{Quersumme $9 + 8 + \dots + 1 = 45$. Alternierend von rechts: $1 - 2 + 3 - 4 + 5 - 6 + 7 - 8 + 9$.}{22pt}
\begin{loesungkarte}{gruen}{1.7}{Quersummen}
\derweg{Quersumme für $9$, alternierende Quersumme für $11$; der Rest der Quersumme ist der Rest der Zahl.}
\schritt{(a)}
Quersumme $45 = 5\cdot9$: durch $9$ teilbar. Alternierende Quersumme (von rechts): $1 - 2 + 3 - 4 + 5 - 6 + 7 - 8 + 9 = 5$: nicht durch $11$ teilbar.
\schritt{(b)}
$n\equiv5\pmod{11}$ (Aufgabe 1.23b: die alternierende Quersumme ist kongruent zur Zahl). Probe: $987\,654\,321 = 11\cdot89\,786\,756 + 5$. \checkmark
\begin{mitnehmen}\kastentext
Quersummenregeln geben nicht nur \enquote{teilbar oder nicht}, sondern den Rest. Für $3$ und $9$: die Quersumme; für $11$: die alternierende; für $7$, $13$: Dreierblöcke alternierend (Aufgabe 1.23c).
\end{mitnehmen}
\end{loesungkarte}

%%%%%%%%%% 1.8
\stupserzwei{tief}{$252 = 1\cdot198 + 54$, $198 = 3\cdot54 + 36$, $54 = 1\cdot36 + 18$, $36 = 2\cdot18$. Rückwärts von $18 = 54 - 36$.}{22pt}
\begin{loesungkarte}{tief}{1.8}{Euklid und Bézout}
\derweg{Vorwärts dividieren bis Rest null; rückwärts jeden Rest durch die Vorgänger ersetzen, bis nur $252$ und $198$ dastehen.}
\schritt{(a)}
$252 = 1\cdot198 + 54$; $198 = 3\cdot54 + 36$; $54 = 1\cdot36 + 18$; $36 = 2\cdot18 + 0$. $\operatorname{ggT} = 18$.
\schritt{(b)}
$18 = 54 - 36 = 54 - (198 - 3\cdot54) = 4\cdot54 - 198 = 4(252 - 198) - 198 = 4\cdot252 - 5\cdot198$. Also $x = 4$, $y = -5$.
\rand{$4\cdot252 = 1008$, $5\cdot198 = 990$, Differenz $18$. \checkmark}{-30pt}
\schritt{(c)}
$12$: nein --- jede Kombination $252x + 198y$ ist ein Vielfaches von $18$, und $18\nmid12$. $36 = 2\cdot18$: ja, $x = 8$, $y = -10$.
\schritt{(d)}
Alle Lösungen von $252x + 198y = 18$: $x = 4 + 11t$, $y = -5 - 14t$, $t\in\mathbb Z$ --- denn $252/18 = 14$ und $198/18 = 11$, und $252\cdot11t - 198\cdot14t = 0$. Probe mit $t = 1$: $15\cdot252 - 19\cdot198 = 3780 - 3762 = 18$. \checkmark
\begin{mitnehmen}\kastentext
Der ggT ist die kleinste positive Zahl der Form $ax + by$, und genau seine Vielfachen sind darstellbar. Rückwärts rechnen ist Fleißarbeit --- die Tabellenform des erweiterten Algorithmus (Band 1, Kapitel 3.4) macht sie mechanisch.
\end{mitnehmen}
\end{loesungkarte}

%%%%%%%%%% 1.9
\stupserzwei{akzent}{$\operatorname{kgV}(4, 6) = 12$, nicht $24$ --- der Faktor $2$ steckt in beiden. Der Beweis in (c) braucht $4x + 6y = 1$, und das gibt es nicht.}{22pt}
\begin{loesungkarte}{akzent}{1.9}{Durch 4 und durch 6}
\derweg{Gemeinsame Teilbarkeit heißt Teilbarkeit durch das kgV; und $\operatorname{kgV} = ab$ nur für teilerfremde $a, b$.}
\schritt{(a)}
$12$: durch $4$ und $6$ teilbar, nicht durch $24$. Ebenso $36$, $60$.
\schritt{(b)}
Durch $\operatorname{kgV}(4, 6) = 12$. Allgemein: $a\mid c$ und $b\mid c$ genau dann, wenn $\operatorname{kgV}(a, b)\mid c$ --- das kgV ist das kleinste gemeinsame Vielfache, und jedes gemeinsame Vielfache ist ein Vielfaches davon (Band 1, Satz 3.7).
\schritt{(c)}
Bézout: $ax + by = 1$ für geeignete $x, y$. Multipliziere mit $c$: $c = acx + bcy$. Nun $b\mid c$, also $ab\mid ac$; und $a\mid c$, also $ab\mid bc$. Beide Summanden sind durch $ab$ teilbar, also $ab\mid c$. $\blacksquare$
\schritt{(d)}
Für $a = 4$, $b = 6$ gibt es keine $x, y$ mit $4x + 6y = 1$ --- die linke Seite ist immer gerade. Der Beweis beginnt mit einer Gleichung, die es nicht gibt. Mit $\operatorname{ggT} = 2$ liefert dieselbe Rechnung nur $4x + 6y = 2$, also $2c = \dots$, also $24\mid2c$, also $12\mid c$ --- genau (b).
\begin{fehler}\kastentext
\enquote{Durch $a$ und durch $b$, also durch $ab$.} Richtig nur für teilerfremde $a, b$. Sonst durch $\operatorname{kgV}(a, b) = ab/\operatorname{ggT}(a, b)$. Die Teilbarkeitsregel für $6$ ($2$ und $3$) funktioniert, die für $8$ als \enquote{$2$ und $4$} nicht.
\end{fehler}
\begin{mitnehmen}\kastentext
Teilerfremdheit ist die Voraussetzung, die man am häufigsten vergisst: beim Zusammensetzen von Teilbarkeiten, beim Kürzen in Kongruenzen (Aufgabe 1.21), beim chinesischen Restsatz (Aufgabe 1.25). Immer dieselbe Zutat: Bézout mit $1$ auf der rechten Seite.
\end{mitnehmen}
\end{loesungkarte}

%%%%%%%%%% 1.10
\stupserzwei{gruen}{$91 = 7\cdot13$, $65 = 5\cdot13$.}{22pt}
\begin{loesungkarte}{gruen}{1.10}{$\operatorname{ggT}(91, 65)$}
\derweg{Drei Divisionen.}
\schritt{Euklid}
$91 = 1\cdot65 + 26$; $65 = 2\cdot26 + 13$; $26 = 2\cdot13 + 0$. $\operatorname{ggT}(91, 65) = 13$.
\probe{$91 = 7\cdot13$, $65 = 5\cdot13$, $\operatorname{ggT}(7, 5) = 1$. \checkmark}
\begin{mitnehmen}\kastentext
Der Algorithmus braucht keine Faktorisierung --- das ist sein Wert bei großen Zahlen, wo Faktorisieren unmöglich ist (Kapitel 19), der ggT aber in Sekunden steht.
\end{mitnehmen}
\end{loesungkarte}

%%%%%%%%%% 1.11
\stupserzwei{gruen}{$13 = 65 - 2\cdot26 = 65 - 2(91 - 65) = 3\cdot65 - 2\cdot91$.}{22pt}
\begin{loesungkarte}{gruen}{1.11}{$91x + 65y = 13$}
\derweg{Euklid rückwärts; dann die allgemeine Lösung mit den gekürzten Koeffizienten.}
\schritt{(a)}
$13 = 65 - 2\cdot26 = 65 - 2\cdot(91 - 65) = 3\cdot65 - 2\cdot91$. Also $x = -2$, $y = 3$: $-182 + 195 = 13$. \checkmark
\schritt{(b)}
$x = -2 + 5t$, $y = 3 - 7t$ ($91/13 = 7$, $65/13 = 5$). Probe $t = 1$: $3\cdot91 - 4\cdot65 = 273 - 260 = 13$. \checkmark
\schritt{(c)}
Jede Zahl $91x + 65y$ ist durch $13$ teilbar; $10$ ist es nicht. Die darstellbaren Zahlen sind genau die Vielfachen des ggT.
\begin{mitnehmen}\kastentext
$ax + by = c$ ist lösbar genau dann, wenn $\operatorname{ggT}(a, b)\mid c$; eine Lösung liefert Euklid rückwärts, alle anderen liegen im Abstand $b/\operatorname{ggT}$ in $x$ und $a/\operatorname{ggT}$ in $y$.
\end{mitnehmen}
\end{loesungkarte}

%%%%%%%%%% 1.12
\stupserzwei{gruen}{$3e = 100 - 5h$ muss durch $3$ teilbar sein: $100 - 5h\equiv1 - 2h\equiv0\pmod3$, also $h\equiv2\pmod3$.}{22pt}
\begin{loesungkarte}{gruen}{1.12}{Hühner und Enten}
\derweg{Eine Lösung finden, alle Lösungen in Schritten $(-3, +5)$ durchlaufen, auf Nichtnegativität prüfen.}
\schritt{(a)}
$5h + 3e = 100$. $h = 20$, $e = 0$ ist eine Lösung; alle: $h = 20 - 3t$, $e = 5t$. Nichtnegativ für $0\le t\le6$: sieben Kombinationen
\[ (h, e)\in\{(20, 0), (17, 5), (14, 10), (11, 15), (8, 20), (5, 25), (2, 30)\} . \]
\schritt{(b)}
$101$: $h\equiv1\pmod3$ (denn $101 - 5h\equiv2 - 2h\equiv0$), $h = 1, 4, \dots, 19$: wieder sieben Lösungen, $(1, 32), (4, 27), \dots, (19, 2)$. $7$: $5h + 3e = 7$ hat keine Lösung in nichtnegativen Zahlen ($h = 0$: $3\nmid7$; $h = 1$: $3\nmid2$).
\schritt{(c)}
Der größte nicht darstellbare Betrag ist $7$: $8 = 5 + 3$, $9 = 3\cdot3$, $10 = 2\cdot5$, und ab drei aufeinanderfolgenden darstellbaren Beträgen ist jeder weitere darstellbar (immer $3$ dazulegen). Allgemein ist $ab - a - b$ der größte nicht darstellbare Betrag für teilerfremde $a, b$ (Sylvester 1882): $15 - 8 = 7$. \checkmark
\begin{center}
\begin{tikzpicture}[x=0.26cm,y=0.26cm,font=\scriptsize]
\draw[grau,-{Latex[length=1.5mm]}] (0,0)--(22,0) node[right]{$h$};
\draw[grau,-{Latex[length=1.5mm]}] (0,0)--(0,34) node[above]{$e$};
\draw[linie] (0,33.3)--(20,0);
\foreach \h/\e in {20/0,17/5,14/10,11/15,8/20,5/25,2/30} {\fill[akzent] (\h,\e) circle (2.2pt);}
\foreach \x in {5,10,15,20} {\node[below,grau] at (\x,0) {\x};}
\foreach \y in {10,20,30} {\node[left,grau] at (0,\y) {\y};}
\end{tikzpicture}
\end{center}
\begin{mitnehmen}\kastentext
Die Lösungen einer linearen diophantischen Gleichung liegen auf einer Geraden im Abstand $(b, -a)/\operatorname{ggT}$; nichtnegative Lösungen sind die Gitterpunkte im ersten Quadranten. Zählen heißt: Zahl der Schritte, die auf der Strecke Platz haben.
\end{mitnehmen}
\end{loesungkarte}

%%%%%%%%%% 1.13
\stupserzwei{gruen}{$\operatorname{kgV} = 91\cdot65/13$.}{22pt}
\begin{loesungkarte}{gruen}{1.13}{$\operatorname{kgV}(91, 65)$}
\derweg{Produkt durch ggT; Begründung über Minimum und Maximum der Exponenten.}
\schritt{(a)}
$\operatorname{kgV}(91, 65) = \frac{91\cdot65}{13} = 7\cdot65 = 455$.
\schritt{(b)}
In der Zerlegung hat $a$ zur Primzahl $p$ den Exponenten $\alpha$, $b$ den Exponenten $\beta$. Der ggT hat $\min(\alpha, \beta)$, das kgV $\max(\alpha, \beta)$, und $\min + \max = \alpha + \beta$ --- der Exponent in $ab$. Das gilt für jede Primzahl, also $\operatorname{ggT}\cdot\operatorname{kgV} = ab$. $\blacksquare$
\probe{$455 = 5\cdot7\cdot13$: durch $91 = 7\cdot13$ und durch $65 = 5\cdot13$ teilbar, und kleiner geht es nicht, weil alle drei Primzahlen gebraucht werden. \checkmark}
\begin{mitnehmen}\kastentext
Für zwei Zahlen gilt $\operatorname{ggT}\cdot\operatorname{kgV} = ab$ --- für drei nicht mehr ($\operatorname{ggT}(2, 2, 2)\cdot\operatorname{kgV}(2, 2, 2) = 4\neq8$). Die Exponentenregel $\min$ und $\max$ gilt dagegen immer.
\end{mitnehmen}
\end{loesungkarte}

%%%%%%%%%% 1.14
\stupserzwei{tief}{$1008 = 16\cdot63 = 2^{4}\cdot3^{2}\cdot7$; $1188 = 4\cdot297 = 2^{2}\cdot3^{3}\cdot11$.}{22pt}
\begin{loesungkarte}{tief}{1.14}{Zwei Zahlen, alle Teiler}
\derweg{Zerlegen, Exponenten vergleichen, Formeln anwenden.}
\schritt{(a)}
$1008 = 2^{4}\cdot3^{2}\cdot7$, $1188 = 2^{2}\cdot3^{3}\cdot11$.
\schritt{(b)}
Teileranzahl $\tau = \prod(\text{Exponent} + 1)$: $\tau(1008) = 5\cdot3\cdot2 = 30$, $\tau(1188) = 3\cdot4\cdot2 = 24$.
\schritt{(c)}
$\operatorname{ggT} = 2^{2}\cdot3^{2} = 36$, $\operatorname{kgV} = 2^{4}\cdot3^{3}\cdot7\cdot11 = 33\,264$. Probe: $36\cdot33\,264 = 1\,197\,504 = 1008\cdot1188$. \checkmark
\schritt{(d)}
Die gemeinsamen Teiler sind die Teiler des ggT: $\tau(36) = 3\cdot3 = 9$ ($1, 2, 3, 4, 6, 9, 12, 18, 36$).
\begin{mitnehmen}\kastentext
Aus der Zerlegung liest man alles ab: $\tau$ aus den Exponenten plus eins, ggT aus den Minima, kgV aus den Maxima, gemeinsame Teiler als Teiler des ggT. Ohne Zerlegung bleibt nur Euklid --- der aber immer geht.
\end{mitnehmen}
\end{loesungkarte}

%%%%%%%%%% 1.15
\stupserzwei{gruen}{$(3 + 1)(2 + 1)(1 + 1) = 24$. Zum Auflisten: alle $2^{a}3^{b}$ für $a\le3$, $b\le2$ aufschreiben, dann jede Zahl noch mit $5$ multiplizieren.}{22pt}
\begin{loesungkarte}{gruen}{1.15}{Die Teiler von 360}
\derweg{Systematisch: erst die Teiler von $72 = 2^{3}\cdot3^{2}$, dann jeden mal $5$.}
\schritt{(a)}
$360 = 2^{3}\cdot3^{2}\cdot5$, $\tau(360) = 4\cdot3\cdot2 = 24$. Teiler von $72$: $1, 2, 4, 8, 3, 6, 12, 24, 9, 18, 36, 72$; mal $5$: $5, 10, 20, 40, 15, 30, 60, 120, 45, 90, 180, 360$. Zusammen $24$. \checkmark
\schritt{(b)}
$240 = 2^{4}\cdot3\cdot5$: $\tau = 5\cdot2\cdot2 = 20$. $288 = 2^{5}\cdot3^{2}$: $6\cdot3 = 18$. $300 = 2^{2}\cdot3\cdot5^{2}$: $3\cdot2\cdot3 = 18$. $336 = 2^{4}\cdot3\cdot7$: $20$. Alle unter $24$. Warum kleine Primzahlen mit fallenden Exponenten: $\tau$ hängt nur von den Exponenten ab, nicht von den Primzahlen --- für dieselben Exponenten ist die Zahl am kleinsten, wenn die größten Exponenten bei den kleinsten Primzahlen stehen ($2^{3}3^{2}5 = 360$ gegen $2\cdot3^{2}5^{3} = 2250$ mit denselben $24$ Teilern).
\begin{mitnehmen}\kastentext
Teiler aufzählen: Gitter aus Exponenten, nicht Probieren. Und $360$, $720$, $5040$ sind \enquote{hochzusammengesetzt} (Ramanujan 1915) --- die Babylonier wählten $360$ für den Kreis, weil man ihn so oft gleichmäßig teilen kann.
\end{mitnehmen}
\end{loesungkarte}

%%%%%%%%%% 1.16
\stupserzwei{gruen}{ggT: je Primzahl der kleinere Exponent, also $2^{2}, 3^{2}, 5^{0}, 7^{0}, 11^{1}$. kgV: der größere.}{22pt}
\begin{loesungkarte}{gruen}{1.16}{ggT und kgV aus der Zerlegung}
\derweg{Je Primzahl Minimum und Maximum der Exponenten.}
\schritt{(a)}
$\operatorname{ggT} = 2^{2}\cdot3^{2}\cdot11 = 396$. $\operatorname{kgV} = 2^{4}\cdot3^{5}\cdot5\cdot7\cdot11^{2} = 16\,465\,680$.
\schritt{(b)}
Exponenten in $ab$: $2^{6}, 3^{7}, 5^{1}, 7^{1}, 11^{3}$; in ggT$\cdot$kgV: $2^{2+4}, 3^{2+5}, 5^{0+1}, 7^{0+1}, 11^{1+2}$. \checkmark
\begin{mitnehmen}\kastentext
Man muss $a$ und $b$ nie ausrechnen. Die Zerlegung ist die bessere Darstellung für alles, was mit Teilbarkeit zu tun hat --- und die schlechtere für Addition, was die Zahlentheorie schwer macht.
\end{mitnehmen}
\end{loesungkarte}

%%%%%%%%%% 1.17
\stupserzwei{gruen}{Nach dem Streichen der Vielfachen von $2$, $3$, $5$, $7$ bleiben genau die Primzahlen --- denn eine zusammengesetzte Zahl $\le50$ hat einen Primfaktor $\le7$.}{22pt}
\begin{loesungkarte}{gruen}{1.17}{Das Sieb bis 50}
\derweg{Vielfache der Primzahlen bis $\sqrt{50}$ streichen; der Rest ist prim.}
\schritt{(a)}
$2, 3, 5, 7, 11, 13, 17, 19, 23, 29, 31, 37, 41, 43, 47$.
\schritt{(b)}
Ist $n\le50$ zusammengesetzt, $n = ab$ mit $2\le a\le b$, so $a^{2}\le ab = n\le50$, also $a\le7$; und $a$ hat einen Primfaktor $\le7$, der $n$ teilt. Jede zusammengesetzte Zahl bis $50$ wird also von $2, 3, 5$ oder $7$ gestrichen.
\schritt{(c)}
$15$ Primzahlen. $50/\ln50 = 50/3{,}912 = 12{,}8$ --- der Primzahlsatz unterschätzt bei kleinen Zahlen; $\operatorname{li}(50)\approx17{,}4$ überschätzt. Beide treffen die Größenordnung (Kapitel 10).
\begin{mitnehmen}\kastentext
Prüfen, ob $n$ prim ist: Teiler bis $\sqrt n$ genügen. Das Sieb macht daraus eine Tabelle; für einzelne große Zahlen ist es nutzlos, und dafür gibt es Kapitel 18.
\end{mitnehmen}
\end{loesungkarte}

%%%%%%%%%% 1.18
\stupserzwei{gruen}{$6k$, $6k + 2$, $6k + 4$ sind gerade; $6k + 3$ ist durch $3$ teilbar. Bleiben $6k + 1$ und $6k + 5 = 6(k + 1) - 1$.}{22pt}
\begin{loesungkarte}{gruen}{1.18}{Primzahlen haben die Form $6k\pm1$}
\derweg{Alle Reste modulo $6$ durchgehen und die ausschließen, die einen Teiler $2$ oder $3$ erzwingen.}
\schritt{(a)}
Jede Zahl ist $6k + r$ mit $r\in\{0, \dots, 5\}$. $r = 0, 2, 4$: gerade; $r = 3$: durch $3$ teilbar. Eine Primzahl $> 3$ ist weder gerade noch durch $3$ teilbar, also $r = 1$ oder $r = 5$: $p = 6k + 1$ oder $p = 6(k + 1) - 1$. $\blacksquare$
\schritt{(b)}
Nein: $25 = 6\cdot4 + 1$ und $35 = 6\cdot6 - 1$ sind nicht prim. Die Form $6k\pm1$ ist notwendig, nicht hinreichend.
\schritt{(c)}
$p^{2} - 1 = (p - 1)(p + 1)$. Mit $p = 6k\pm1$: $(p - 1)(p + 1) = 6k(6k\pm2) = 12k(3k\pm1)$. Von $k$ und $3k\pm1$ ist eines gerade (ist $k$ ungerade, so ist $3k\pm1$ gerade), also $24\mid p^{2} - 1$. $\blacksquare$ Probe: $p = 5$: $24$; $p = 7$: $48$; $p = 11$: $120 = 5\cdot24$. \checkmark
\begin{mitnehmen}\kastentext
\enquote{Primzahlen $> 3$ sind $\equiv\pm1\pmod6$} sagt: Von sechs aufeinanderfolgenden Zahlen kommen nur zwei als Primzahlen infrage. Ein Primzahltest muss nur ein Drittel aller Kandidaten prüfen --- die \enquote{Wheel-Faktorisierung} der Praxis.
\end{mitnehmen}
\end{loesungkarte}

%%%%%%%%%% 1.19
\stupserzwei{gruen}{$N = 4p_1\cdots p_k - 1\equiv3\pmod4$. Wären alle Primfaktoren von $N$ $\equiv1\pmod4$, wäre $N$ es auch (nach (a)). Also hat $N$ einen Primfaktor $\equiv3$, und der ist keines der $p_i$.}{22pt}
\begin{loesungkarte}{gruen}{1.19}{Unendlich viele $4k + 3$}
\derweg{Euklid: Aus endlich vielen eine neue Zahl bauen, die einen neuen Primfaktor der gewünschten Form haben muss.}
\schritt{(a)}
$(4a + 1)(4b + 1) = 16ab + 4a + 4b + 1 = 4(4ab + a + b) + 1$. Induktiv für beliebig viele Faktoren. $\blacksquare$
\schritt{(b)}
Angenommen, $p_1, \dots, p_k$ seien alle Primzahlen der Form $4k + 3$. Setze $N = 4p_1p_2\cdots p_k - 1$. $N$ ist ungerade und $\equiv3\pmod4$. Alle Primfaktoren von $N$ sind ungerade, also $\equiv1$ oder $\equiv3\pmod4$. Wären alle $\equiv1$, so wäre $N\equiv1$ nach (a) --- Widerspruch. Also hat $N$ einen Primfaktor $q\equiv3\pmod4$. Aber kein $p_i$ teilt $N$ (es ließe Rest $-1$, also $p_i - 1$). Also ist $q$ eine neue Primzahl der Form $4k + 3$. $\blacksquare$
\schritt{(c)}
Für $4k + 1$ scheitert das Analogon: Ein Produkt von Zahlen $\equiv3\pmod4$ kann $\equiv1$ sein ($3\cdot3 = 9$), also folgt aus $N\equiv1$ nicht, dass ein Primfaktor $\equiv1$ ist. Für $4k + 1$ braucht man $N = (2p_1\cdots p_k)^{2} + 1$ und den Satz, dass Primteiler von $x^{2} + 1$ immer $\equiv1\pmod4$ sind (Ergänzungssatz, Kapitel 8). Für allgemeines $ak + b$ mit $\operatorname{ggT}(a, b) = 1$: Dirichlet 1837, analytisch (Kapitel 10).
\begin{mitnehmen}\kastentext
Euklids Beweis ist ein Bauplan: Konstruiere aus den bekannten Primzahlen eine Zahl, die eine unbekannte erzwingt. Für Restklassen funktioniert er genau dann elementar, wenn die \enquote{anderen} Restklassen multiplikativ abgeschlossen sind.
\end{mitnehmen}
\end{loesungkarte}

%%%%%%%%%% 1.20
\stupserzwei{tief}{$14x\equiv30\pmod{100}\iff7x\equiv15\pmod{50}$. Inverses von $7$ modulo $50$: $7\cdot43 = 301 = 6\cdot50 + 1$.}{22pt}
\begin{loesungkarte}{tief}{1.20}{Eine lineare Kongruenz mit zwei Lösungen}
\derweg{ggT prüfen, durch ihn kürzen (auch den Modul), dann mit dem Inversen multiplizieren; am Ende die Lösungen modulo des alten Moduls auffächern.}
\schritt{(a)}
$\operatorname{ggT}(14, 100) = 2\mid30$: lösbar, zwei Lösungen. Kürzen: $7x\equiv15\pmod{50}$. $7^{-1}\equiv43\pmod{50}$, also $x\equiv15\cdot43 = 645\equiv45\pmod{50}$. Modulo $100$: $x\equiv45$ oder $x\equiv95$. Probe: $14\cdot45 = 630\equiv30$, $14\cdot95 = 1330\equiv30$. \checkmark
\schritt{(b)}
$2\nmid31$: keine Lösung --- $14x$ ist immer gerade, $31 + 100k$ immer ungerade.
\schritt{(c)}
Erweiterter Euklid für $7$ und $50$: $50 = 7\cdot7 + 1$, also $1 = 50 - 7\cdot7$, also $7\cdot(-7)\equiv1$, $7^{-1}\equiv-7\equiv43$. Damit $x\equiv15\cdot43\equiv45\pmod{50}$, wie in (a).
\begin{mitnehmen}\kastentext
$ax\equiv b\pmod m$: $d = \operatorname{ggT}(a, m)$ muss $b$ teilen; dann durch $d$ kürzen, Inverses nehmen, und die eine Lösung modulo $m/d$ wird zu $d$ Lösungen modulo $m$. Das Inverse liefert immer Euklid --- oder, bei kleinen Zahlen, das Hinsehen.
\end{mitnehmen}
\end{loesungkarte}

%%%%%%%%%% 1.21
\stupserzwei{akzent}{Aus $ca\equiv cb\pmod m$ folgt $m\mid c(a - b)$. Teilt man $m$ und $c$ durch ihren ggT $d$, so ist $m/d$ teilerfremd zu $c/d$ und muss $a - b$ teilen.}{22pt}
\begin{loesungkarte}{akzent}{1.21}{Kürzen in Kongruenzen}
\derweg{Kongruenz als Teilbarkeit lesen; der Modul verliert beim Kürzen den Faktor, den er mit $c$ gemeinsam hat.}
\schritt{(a)}
$6x\bmod9$ für $x = 0, \dots, 8$: $0, 6, 3, 0, 6, 3, 0, 6, 3$. Gleich $6$ für $x = 1, 4, 7$ --- drei Lösungen, nicht eine.
\schritt{(b)}
Aus $ca\equiv cb\pmod m$ folgt $a\equiv b\pmod{m/\operatorname{ggT}(c, m)}$. Beweis: $m\mid c(a - b)$; mit $d = \operatorname{ggT}(c, m)$ ist $\frac md\mid\frac cd(a - b)$ und $\operatorname{ggT}(\frac md, \frac cd) = 1$, also $\frac md\mid a - b$ (Lemma von Euklid). Hier: $6x\equiv6\pmod9$ gibt $x\equiv1\pmod3$, also $x\in\{1, 4, 7\}$ modulo $9$. \checkmark
\schritt{(c)}
Genau für $\operatorname{ggT}(c, m) = 1$ bleibt der Modul: Dann hat $c$ ein Inverses modulo $m$, und Kürzen ist Multiplikation mit $c^{-1}$ --- eine umkehrbare Operation. Ist $\operatorname{ggT}(c, m) > 1$, hat $c$ kein Inverses, und $c\cdot(\cdot)$ ist nicht injektiv modulo $m$: Verschiedene $x$ geben dasselbe $cx$.
\schritt{(d)}
$15x\equiv25\pmod{35}$: $\operatorname{ggT}(15, 35) = 5$, kürzen zu $3x\equiv5\pmod7$. $3^{-1}\equiv5\pmod7$ ($15\equiv1$), also $x\equiv25\equiv4\pmod7$: $x\in\{4, 11, 18, 25, 32\}$ modulo $35$. Probe: $15\cdot4 = 60\equiv25$. \checkmark
\begin{fehler}\kastentext
\enquote{$6x\equiv6$, also $x\equiv1$.} --- Man hat durch $6$ geteilt, das modulo $9$ kein Inverses hat. Die Restklassen sind kein Körper, wenn der Modul nicht prim ist: Nullteiler ($3\cdot6\equiv0$) machen die Division zweideutig.
\end{fehler}
\begin{mitnehmen}\kastentext
Kürzen durch $c$: Modul durch $\operatorname{ggT}(c, m)$ teilen. Ist $m$ prim, darf man alles außer $0$ kürzen --- der Grund, warum $\mathbb Z/p\mathbb Z$ ein Körper ist und $\mathbb Z/9\mathbb Z$ nicht.
\end{mitnehmen}
\end{loesungkarte}

%%%%%%%%%% 1.22
\stupserzwei{gruen}{$7^{100}\equiv2^{100} = (2^{4})^{25}\equiv1^{25}$.}{22pt}
\begin{loesungkarte}{gruen}{1.22}{$7^{100}$ modulo 5}
\derweg{Basis reduzieren, Periode finden, Exponent modulo Periode.}
\schritt{(a)}
$7\equiv2\pmod5$; $2^{1}, 2^{2}, 2^{3}, 2^{4}\equiv2, 4, 3, 1$. Periode $4$, $100 = 4\cdot25$: $7^{100}\equiv1$. Rest $1$.
\schritt{(b)}
$7^{101}\equiv2^{1} = 2$. $7^{2026}$: $2026 = 4\cdot506 + 2$, also $\equiv2^{2} = 4$.
\probe{$7^{2} = 49\equiv4$, $7^{4}\equiv16\equiv1$. \checkmark}
\begin{mitnehmen}\kastentext
Die Potenzen einer Zahl modulo $m$ sind periodisch; die Periode heißt Ordnung und teilt $\varphi(m)$ (Kapitel 7). Exponenten rechnet man modulo der Ordnung --- nie modulo $m$.
\end{mitnehmen}
\end{loesungkarte}

%%%%%%%%%% 1.23
\stupserzwei{gruen}{$n = \sum a_k10^{k}\equiv\sum a_k\cdot1^{k} = Q(n)\pmod9$. Für $11$: $10\equiv-1$, also $10^{k}\equiv(-1)^{k}$.}{22pt}
\begin{loesungkarte}{gruen}{1.23}{Die Neunerprobe}
\derweg{Die Zahl als Summe $a_k10^{k}$ schreiben und $10$ durch seinen Rest ersetzen.}
\schritt{(a)}
$10\equiv1\pmod9$, also $10^{k}\equiv1$ für alle $k$, also $n = \sum_ka_k10^{k}\equiv\sum_ka_k = Q(n)\pmod9$. Insbesondere $9\mid n\iff9\mid Q(n)$ --- und mehr: $n$ und $Q(n)$ lassen denselben Rest. (Dasselbe für $3$.) $\blacksquare$
\schritt{(b)}
$10\equiv-1\pmod{11}$, also $10^{k}\equiv(-1)^{k}$ und $n\equiv a_0 - a_1 + a_2 - \dots = Q'(n)$, die alternierende Quersumme von rechts. $\blacksquare$
\schritt{(c)}
$10^{3} = 1000 = 143\cdot7 - 1\equiv-1\pmod7$. Zerlegt man $n$ in Dreierblöcke von rechts, $n = B_0 + B_1\cdot10^{3} + B_2\cdot10^{6} + \dots$, so $n\equiv B_0 - B_1 + B_2 - \dots\pmod7$. Beispiel $987\,654\,321$: $321 - 654 + 987 = 654 = 7\cdot93 + 3$, also Rest $3$. (Dieselbe Regel gilt für $11$ und $13$, denn $1001 = 7\cdot11\cdot13$.)
\probe{(c): $987\,654\,321 = 7\cdot141\,093\,474 + 3$. \checkmark}
\begin{mitnehmen}\kastentext
Jede Teilbarkeitsregel ist \enquote{$10^{k}$ modulo $m$ ausgerechnet}. Für $m = 9$: alles $1$. Für $11$: $\pm1$. Für $7$: Periode $6$, oder Dreierblöcke mit $\pm1$. Für $2$ und $5$: nur $10^{0}$ zählt --- die letzte Ziffer.
\end{mitnehmen}
\end{loesungkarte}

%%%%%%%%%% 1.24
\stupserzwei{gruen}{$2x\equiv5\pmod7$: $2^{-1}\equiv4$, also $x\equiv20\equiv6\pmod7$. Zurück modulo $21$: $6, 13, 20$.}{22pt}
\begin{loesungkarte}{gruen}{1.24}{$6x\equiv15\pmod{21}$}
\derweg{Durch $\operatorname{ggT}(6, 21) = 3$ kürzen, lösen, auffächern.}
\schritt{Rechnung}
$3\mid15$: lösbar, drei Lösungen. $2x\equiv5\pmod7$; $2\cdot4 = 8\equiv1$, also $x\equiv4\cdot5 = 20\equiv6\pmod7$. Modulo $21$: $x\in\{6, 13, 20\}$.
\probe{$6\cdot6 = 36\equiv15$; $6\cdot13 = 78 = 3\cdot21 + 15$; $6\cdot20 = 120 = 5\cdot21 + 15$. \checkmark}
\begin{mitnehmen}\kastentext
Drei Lösungen, weil $\operatorname{ggT} = 3$; eine Lösung modulo $7$, drei modulo $21$. Die Zahl der Lösungen steht fest, bevor man rechnet.
\end{mitnehmen}
\end{loesungkarte}

%%%%%%%%%% 1.25
\stupserzwei{tief}{$e_1 = 35\cdot3 = 105$ ($35\equiv3\pmod4$, $3\cdot3\equiv1$), $e_2 = 28\cdot2 = 56$, $e_3 = 20\cdot6 = 120$. Dann $x = 1\cdot e_1 + 2\cdot e_2 + 3\cdot e_3$.}{22pt}
\begin{loesungkarte}{tief}{1.25}{Drei Reste, ein Rest}
\derweg{Für jeden Modul einen Baustein $e_i$ bauen, der dort $1$ und sonst $0$ ist; die Lösung ist die gewichtete Summe.}
\schritt{(a), (b)}
$M = 4\cdot5\cdot7 = 140$. $e_1$: Vielfaches von $35$ mit Rest $1$ bei $4$: $35\equiv3\pmod4$, $3^{-1}\equiv3$, $e_1 = 35\cdot3 = 105$. $e_2$: Vielfaches von $28$ mit Rest $1$ bei $5$: $28\equiv3\pmod5$, $3^{-1}\equiv2$, $e_2 = 56$. $e_3$: Vielfaches von $20$ mit Rest $1$ bei $7$: $20\equiv6\equiv-1\pmod7$, $(-1)^{-1} = -1\equiv6$, $e_3 = 120$.
\[ x\equiv1\cdot105 + 2\cdot56 + 3\cdot120 = 577\equiv577 - 4\cdot140 = 17\pmod{140} . \]
Zweites System: $x\equiv3\cdot105 + 0\cdot56 + 6\cdot120 = 1035\equiv1035 - 7\cdot140 = 55\pmod{140}$.
\rand{$17$: Rest $1$ bei $4$, $2$ bei $5$, $3$ bei $7$. \checkmark\ $55$: $3$, $0$, $6$. \checkmark}{-40pt}
\schritt{(c)}
Die Bausteine brauchen $M/m_i$ invertierbar modulo $m_i$, also $\operatorname{ggT}(M/m_i, m_i) = 1$ --- das ist paarweise Teilerfremdheit. Bei $4$ und $6$: $x\equiv1\pmod4$ heißt $x$ ungerade, $x\equiv2\pmod6$ heißt $x$ gerade. Keine Lösung. Und $x\equiv1\pmod4$, $x\equiv3\pmod6$ hätte Lösungen ($9, 21, \dots$), aber modulo $\operatorname{kgV} = 12$, nicht modulo $24$. Allgemein: lösbar genau dann, wenn die Reste modulo $\operatorname{ggT}(m_i, m_j)$ übereinstimmen.
\begin{mitnehmen}\kastentext
Die Bausteine $e_i$ sind die \enquote{Einheitsvektoren} des Restsatzes: Mit ihnen ist jedes System zu denselben Moduln eine Zeile. Sie sind genau das, was $\mathbb Z/140\cong\mathbb Z/4\times\mathbb Z/5\times\mathbb Z/7$ bedeutet.
\end{mitnehmen}
\end{loesungkarte}

%%%%%%%%%% 1.26
\stupserzwei{gruen}{Rest $2$ bei $3$ und $7$: $x\equiv2\pmod{21}$, also $x\in\{2, 23, 44, \dots\}$. Welche davon lässt Rest $3$ bei $5$?}{22pt}
\begin{loesungkarte}{gruen}{1.26}{Meister Sun}
\derweg{Zwei Kongruenzen zusammenfassen, die dritte durch Probieren in Schritten des Produkts.}
\schritt{(a)}
Aus $x\equiv2\pmod3$ und $x\equiv2\pmod7$ folgt $x\equiv2\pmod{21}$: $2, 23, 44, 65, 86$. Rest $3$ bei $5$: $23$. Also $x = 23$.
\schritt{(b)}
$x\equiv23\pmod{105}$: $23, 128, 233, \dots$
\probe{$23 = 7\cdot3 + 2 = 4\cdot5 + 3 = 3\cdot7 + 2$. \checkmark}
\begin{mitnehmen}\kastentext
Meister Sun selbst rechnete mit den Bausteinen $70, 21, 15$ (das sind $e_1, e_2, e_3$ für $3, 5, 7$): $2\cdot70 + 3\cdot21 + 2\cdot15 = 233$, minus $2\cdot105$ gleich $23$. Der Vers dazu steht in Band 1, Kapitel 6.
\end{mitnehmen}
\end{loesungkarte}

%%%%%%%%%% 1.27
\stupserzwei{gruen}{Modulo $3$, $5$, $7$ je $x\equiv\pm1$: $2\cdot2\cdot2 = 8$ Kombinationen, jede nach dem Restsatz genau eine Zahl modulo $105$.}{22pt}
\begin{loesungkarte}{gruen}{1.27}{$x^{2}\equiv1\pmod{105}$}
\derweg{Modulo jeder Primzahl zwei Lösungen, der Restsatz kombiniert sie unabhängig.}
\schritt{(a)}
$x^{2}\equiv1\pmod p$ hat für ungerades $p$ genau die Lösungen $\pm1$ (denn $p\mid(x - 1)(x + 1)$ erzwingt $p\mid x - 1$ oder $p\mid x + 1$). Modulo $105 = 3\cdot5\cdot7$: $2^{3} = 8$ Vorzeichenkombinationen, je eine Lösung. Mit den Bausteinen $70, 21, 15$: $x\equiv\pm70\pm21\pm15$. Die acht Lösungen: $1, 29, 34, 41, 64, 71, 76, 104$.
\rand{$29^{2} = 841 = 8\cdot105 + 1$. \checkmark\ $34^{2} = 1156 = 11\cdot105 + 1$. \checkmark}{-30pt}
\schritt{(b)}
$2^{k}$ Lösungen --- jede Primzahl verdoppelt.
\schritt{(c)}
$x = 29$: $\operatorname{ggT}(28, 105) = 7$. $x = 34$: $\operatorname{ggT}(33, 105) = 3$. $x = 41$: $\operatorname{ggT}(40, 105) = 5$. Jede nichttriviale Quadratwurzel von $1$ liefert einen echten Faktor von $105$ --- denn $105\mid(x - 1)(x + 1)$, aber weder $x - 1$ noch $x + 1$: Die Primfaktoren verteilen sich auf beide.
\begin{mitnehmen}\kastentext
Nichttriviale Quadratwurzeln von $1$ gibt es nur modulo zusammengesetzter Zahlen, und jede verrät einen Faktor. Das ist der Kern des Miller--Rabin-Tests (Kapitel 18) und der Faktorisierung mit Differenzen von Quadraten (Kapitel 19).
\end{mitnehmen}
\end{loesungkarte}

%%%%%%%%%% 1.28
\stupserzwei{gruen}{$x\equiv1\pmod{60}$ und $x\equiv0\pmod7$: $x = 60k + 1$, und $60k + 1\equiv4k + 1\equiv0\pmod7$ gibt $k\equiv5$.}{22pt}
\begin{loesungkarte}{gruen}{1.28}{Die Eier im Korb}
\derweg{Die fünf Bedingungen zu einer zusammenfassen (kgV), dann mit der sechsten kombinieren.}
\schritt{Rechnung}
Rest $1$ bei $2, 3, 4, 5, 6$ heißt $x - 1$ durch $\operatorname{kgV}(2, 3, 4, 5, 6) = 60$ teilbar: $x = 60k + 1$. Dazu $7\mid x$: $60k + 1\equiv4k + 1\equiv0\pmod7$, $4k\equiv6$, $k\equiv6\cdot4^{-1} = 6\cdot2 = 12\equiv5\pmod7$. Also $x = 60\cdot5 + 1 = 301$; alle Lösungen $301 + 420t$.
\probe{$301 = 7\cdot43$; $300 = 60\cdot5$ ist durch $2, 3, 4, 5, 6$ teilbar. \checkmark}
\begin{mitnehmen}\kastentext
Moduln, die nicht teilerfremd sind, fasst man erst zusammen (kgV, Verträglichkeit der Reste prüfen), dann wendet man den Restsatz an. Hier passten alle fünf Reste zusammen, weil sie gleich waren.
\end{mitnehmen}
\end{loesungkarte}

%%%%%%%%%% 1.29
\stupserzwei{tief}{$3^{1}, 3^{2}, 3^{3}, 3^{4}\equiv3, 9, 7, 1\pmod{10}$; $2025 = 4\cdot506 + 1$. Für (b): $1000 = 12\cdot83 + 4$.}{22pt}
\begin{loesungkarte}{tief}{1.29}{Letzte Ziffern und Ordnungen}
\derweg{Ordnung oder Fermat liefert die Periode; der Exponent wird modulo Periode reduziert.}
\schritt{(a)}
Modulo $10$: $3^{4} = 81\equiv1$, Periode $4$. $2025\equiv1\pmod4$, also $3^{2025}\equiv3^{1} = 3$: letzte Ziffer $3$.
\schritt{(b)}
Fermat: $2^{12}\equiv1\pmod{13}$. $1000 = 12\cdot83 + 4$, also $2^{1000}\equiv2^{4} = 16\equiv3\pmod{13}$.
\schritt{(c)}
Potenzen von $2$ modulo $13$: $2, 4, 8, 3, 6, 12, 11, 9, 5, 10, 7, 1$ --- erst $2^{12}$ ist $1$: Ordnung $12$, $2$ ist Primitivwurzel. Eine kleinere Ordnung hätte den Exponenten weiter reduziert; hier war $12$ schon das Beste. (Für $3$ modulo $13$: $3^{3} = 27\equiv1$, Ordnung $3$ --- da reduziert man $1000$ modulo $3$.)
\begin{mitnehmen}\kastentext
Der Exponent lebt modulo der Ordnung; die Ordnung teilt $\varphi(m)$; Fermat/Euler geben $\varphi(m)$ als sichere, oft zu große Periode. Wer die wahre Ordnung kennt, rechnet mit kleineren Zahlen.
\end{mitnehmen}
\end{loesungkarte}

%%%%%%%%%% 1.30
\stupserzwei{akzent}{$2^{10} = 1024 = 93\cdot11 + 1\equiv1\pmod{11}$ und $1024 = 33\cdot31 + 1\equiv1\pmod{31}$. Also $2^{340} = (2^{10})^{34}\equiv1$ modulo beider Faktoren.}{22pt}
\begin{loesungkarte}{akzent}{1.30}{341 besteht den Test}
\derweg{Modulo beider Primfaktoren rechnen und mit dem Restsatz zusammensetzen; dann sehen, dass eine andere Basis die Zahl entlarvt.}
\schritt{(a)}
$2^{10} = 1024\equiv1\pmod{11}$ und $\equiv1\pmod{31}$. Also $2^{340} = (2^{10})^{34}\equiv1$ modulo $11$ und $31$, also modulo $341$. Die zusammengesetzte Zahl $341$ \enquote{besteht} den Fermat-Test zur Basis $2$: eine \emph{Pseudoprimzahl} zur Basis $2$, die kleinste.
\schritt{(b)}
$3^{5} = 243\equiv1\pmod{11}$ ($243 = 22\cdot11 + 1$), also $3^{340} = (3^{5})^{68}\equiv1\pmod{11}$. Modulo $31$: $3^{3} = 27\equiv-4$, $3^{5} = 243 = 7\cdot31 + 26\equiv-5$, $3^{10}\equiv25\equiv-6$, $3^{30}\equiv(-6)^{3} = -216 = -7\cdot31 + 1\equiv1$; also $3^{340} = 3^{330}\cdot3^{10}\equiv1\cdot(-6)\equiv25\pmod{31}$. Zusammen: $x\equiv1\pmod{11}$, $x\equiv25\pmod{31}$ --- das ist $56$ ($56 = 5\cdot11 + 1$, $56 = 31 + 25$). $3^{340}\equiv56\not\equiv1\pmod{341}$: Zur Basis $3$ fällt $341$ durch. Der Test ist zusammengesetzt, aber die Basis $2$ hat es nicht gemerkt.
\schritt{(c)}
Fermat: \emph{Wenn} $n$ prim ist, \emph{dann} $a^{n-1}\equiv1$ für alle $a$ mit $p\nmid a$. Gilt $a^{n-1}\not\equiv1$ für ein $a$, ist $n$ sicher zusammengesetzt. Gilt $a^{n-1}\equiv1$, ist $n$ prim \emph{oder} eine Pseudoprimzahl zur Basis $a$ --- der Test beweist nie Primalität.
\schritt{(d)}
Weil Pseudoprimzahlen selten sind (unter $10^{10}$ etwa $15\,000$ zur Basis $2$ gegen $455$ Millionen Primzahlen) und weil eine zweite Basis die meisten entlarvt. Nur die Carmichael-Zahlen ($561, 1105, 1729, \dots$) bestehen den Test für \emph{jede} teilerfremde Basis --- und für die gibt es Miller--Rabin (Kapitel 18, Aufgabe 18.1).
\begin{fehler}\kastentext
\enquote{$a^{n-1}\equiv1\pmod n$, also ist $n$ prim.} --- Das ist die Umkehrung des kleinen Fermat, und sie ist falsch: $341$, $561$, $645$ bestehen den Test zur Basis $2$. Ein Satz der Form \enquote{prim $\Rightarrow$ Eigenschaft} ist ein Test, der nur \emph{aussortieren} kann.
\end{fehler}
\begin{mitnehmen}\kastentext
Implikation ist keine Äquivalenz. Fermat liefert einen \emph{Zusammengesetztheitsbeweis}, nie einen Primzahlbeweis; dass er in der Praxis trotzdem reicht, liegt an Wahrscheinlichkeiten, nicht an Logik.
\end{mitnehmen}
\end{loesungkarte}

%%%%%%%%%% 1.31
\stupserzwei{gruen}{$3^{1000} = (3^{6})^{166}\cdot3^{4}$.}{22pt}
\begin{loesungkarte}{gruen}{1.31}{$3^{1000}$ modulo 7}
\derweg{Fermat, Exponent modulo $6$.}
\schritt{Rechnung}
$3^{6}\equiv1\pmod7$; $1000 = 6\cdot166 + 4$; $3^{1000}\equiv3^{4} = 81 = 11\cdot7 + 4\equiv4$.
\probe{Ordnung von $3$ modulo $7$ ist $6$ ($3, 2, 6, 4, 5, 1$), also war Fermat hier scharf. \checkmark}
\begin{mitnehmen}\kastentext
Eine Zeile --- wenn man den Exponenten reduziert, bevor man rechnet.
\end{mitnehmen}
\end{loesungkarte}

%%%%%%%%%% 1.32
\stupserzwei{gruen}{$\varphi(2^{3}) = 4$, $\varphi(3^{2}) = 6$, $\varphi(5) = 4$; $\varphi(1000) = \varphi(2^{3})\varphi(5^{3}) = 4\cdot100$.}{22pt}
\begin{loesungkarte}{gruen}{1.32}{$\varphi(360)$ und $\varphi(1000)$}
\derweg{Zerlegen, $\varphi(p^{k}) = p^{k-1}(p - 1)$, multiplizieren.}
\schritt{(a)}
$\varphi(360) = \varphi(8)\varphi(9)\varphi(5) = 4\cdot6\cdot4 = 96$. $\varphi(1000) = \varphi(8)\varphi(125) = 4\cdot100 = 400$.
\schritt{(b)}
$360\cdot(1 - \tfrac12)(1 - \tfrac13)(1 - \tfrac15) = 360\cdot\tfrac12\cdot\tfrac23\cdot\tfrac45 = 96$. \checkmark
\schritt{(c)}
$\varphi(n) = n - 1$ genau für Primzahlen: Dann sind alle $1, \dots, n - 1$ teilerfremd zu $n$; ist $n$ zusammengesetzt, ist ein echter Teiler nicht teilerfremd, also $\varphi(n)\le n - 2$.
\begin{mitnehmen}\kastentext
$\varphi(n)$ zählt die invertierbaren Reste. Es ist leicht aus der Zerlegung und --- für große $n$ --- ohne sie nicht zu berechnen; und wer $\varphi(n)$ kennt, kann $n = pq$ faktorisieren (Kapitel 17). Darauf beruht RSA.
\end{mitnehmen}
\end{loesungkarte}

%%%%%%%%%% 1.33
\stupserzwei{gruen}{$7^{4} = 2401\equiv1\pmod{100}$. $2026 = 4\cdot506 + 2$.}{22pt}
\begin{loesungkarte}{gruen}{1.33}{Die letzten beiden Ziffern von $7^{2026}$}
\derweg{Die Ordnung von $7$ modulo $100$ finden --- sie ist viel kleiner als $\varphi(100) = 40$.}
\schritt{(a)}
$7^{2} = 49$, $7^{4} = 49^{2} = 2401\equiv1\pmod{100}$. Also $7^{2026} = (7^{4})^{506}\cdot7^{2}\equiv49$. Die letzten Ziffern sind $49$.
\schritt{(b)}
Ordnung $4$. Nach Euler teilt die Ordnung $\varphi(100) = 40$; möglich wären $1, 2, 4, 5, 8, 10, 20, 40$. Mit Euler allein: $2026 = 40\cdot50 + 26$, $7^{26} = (7^{4})^{6}\cdot7^{2}\equiv49$ --- dasselbe, nur mit einem Umweg.
\begin{mitnehmen}\kastentext
Euler gibt eine Periode, die Ordnung die kleinste. Bei kleinen Basen lohnt es, ein paar Potenzen auszurechnen, bevor man $\varphi$ bemüht: $7^{4}\equiv1$ ist schneller gefunden als $7^{40}$ ausgerechnet.
\end{mitnehmen}
\end{loesungkarte}

%%%%%%%%%% 1.34
\stupserzwei{gruen}{$2^{3} = 8\equiv1$: Ordnung $3$. $3^{1}, \dots, 3^{6}\equiv3, 2, 6, 4, 5, 1$: Ordnung $6$. $6 = -1$: Ordnung $2$.}{22pt}
\begin{loesungkarte}{gruen}{1.34}{Ordnungen modulo 7}
\derweg{Potenzen bilden bis zur $1$; die Ordnungen sind Teiler von $6$.}
\schritt{(a)}
$\operatorname{ord}(1) = 1$; $\operatorname{ord}(2) = 3$ ($2, 4, 1$); $\operatorname{ord}(3) = 6$ ($3, 2, 6, 4, 5, 1$); $\operatorname{ord}(4) = 3$ ($4, 2, 1$); $\operatorname{ord}(5) = 6$ ($5, 4, 6, 2, 3, 1$); $\operatorname{ord}(6) = 2$.
\schritt{(b)}
Primitivwurzeln: $3$ und $5$ (Ordnung $6 = \varphi(7)$). Formel $\varphi(\varphi(7)) = \varphi(6) = 2$. \checkmark
\schritt{(c)}
Ordnungen $1, 2, 3, 3, 6, 6$ --- alle Teiler von $6$, und jeder Teiler $d$ kommt genau $\varphi(d)$-mal vor: $\varphi(1) = 1$, $\varphi(2) = 1$, $\varphi(3) = 2$, $\varphi(6) = 2$. Das ist kein Zufall, sondern die Struktur der zyklischen Gruppe (Band 11).
\begin{mitnehmen}\kastentext
Die Ordnung teilt $\varphi(m)$ (Lagrange); die Zahl der Elemente der Ordnung $d$ ist $\varphi(d)$, wenn eine Primitivwurzel existiert. Modulo einer Primzahl ist $(\mathbb Z/p)^{\times}$ zyklisch --- der Satz, den Kapitel 7.5 beweist.
\end{mitnehmen}
\end{loesungkarte}

%%%%%%%%%% 1.35
\stupserzwei{gruen}{$g^{4}\not\equiv1$ und $g^{6}\not\equiv1$ prüfen. $2^{4} = 16\equiv3$, $2^{6} = 64\equiv12$: $2$ ist Primitivwurzel.}{22pt}
\begin{loesungkarte}{gruen}{1.35}{Primitivwurzeln modulo 13}
\derweg{Ordnung $12$ heißt: keine der Potenzen $g^{12/q}$ für $q\in\{2, 3\}$ ist $1$.}
\schritt{(a)}
Die Ordnung teilt $12$; sie ist genau $12$, wenn $g^{6}\not\equiv1$ und $g^{4}\not\equiv1$. Tabelle: $g^{2}, g^{4} = (g^{2})^{2}, g^{6} = g^{4}g^{2}$ modulo $13$:
\begin{center}
{\small\renewcommand{\arraystretch}{1.25}
\begin{tabular}{@{}c|cccccccccccc@{}}
$g$ & 1 & 2 & 3 & 4 & 5 & 6 & 7 & 8 & 9 & 10 & 11 & 12\\\hline
$g^{4}$ & 1 & 3 & 3 & 9 & 1 & 9 & 9 & 1 & 9 & 3 & 3 & 1\\
$g^{6}$ & 1 & 12 & 1 & 1 & 12 & 12 & 12 & 12 & 1 & 1 & 12 & 1\\
\end{tabular}}
\end{center}
Beide $\neq1$ bei $g = 2, 6, 7, 11$. Das sind die Primitivwurzeln.
\schritt{(b)}
$\varphi(12) = 4$. \checkmark\ (Es sind $2^{k}$ für $k$ teilerfremd zu $12$: $2^{1} = 2$, $2^{5} = 6$, $2^{7} = 11$, $2^{11} = 7$.)
\schritt{(c)}
$2^{0}, \dots, 2^{11}$: $1, 2, 4, 8, 3, 6, 12, 11, 9, 5, 10, 7$. $5 = 2^{9}$, also $\log_25 = 9$. Alle Reste $1, \dots, 12$ kommen vor --- das ist, was Primitivwurzel heißt.
\begin{mitnehmen}\kastentext
Primitivwurzel-Test: $g^{(p-1)/q}\not\equiv1$ für jeden Primteiler $q$ von $p - 1$. Der diskrete Logarithmus ist leicht abzulesen, wenn die Tabelle da ist, und ohne Tabelle schwer --- Diffie--Hellman (Kapitel 17) lebt davon.
\end{mitnehmen}
\end{loesungkarte}

%%%%%%%%%% 1.36
\stupserzwei{gruen}{Modulo $11$: $2\cdot6 = 12\equiv1$, $3\cdot4\equiv1$, $5\cdot9 = 45\equiv1$, $7\cdot8 = 56\equiv1$. Übrig bleiben $1$ und $10$.}{22pt}
\begin{loesungkarte}{gruen}{1.36}{Wilson für 11 und 9}
\derweg{Jedes Element außer $\pm1$ mit seinem Inversen paaren; die Paare geben $1$, übrig bleibt $-1$.}
\schritt{(a)}
$10! = 1\cdot(2\cdot6)(3\cdot4)(5\cdot9)(7\cdot8)\cdot10\equiv1\cdot1\cdot1\cdot1\cdot1\cdot10\equiv-1\pmod{11}$. \checkmark\ Nur $1$ und $10\equiv-1$ sind zu sich selbst invers ($x^{2}\equiv1$ hat modulo einer Primzahl nur $\pm1$, Aufgabe 1.27).
\schritt{(b)}
$8! = 40320 = 4480\cdot9$, also $8!\equiv0\not\equiv-1\pmod9$. Nach Wilson (Richtung \enquote{$(n - 1)!\equiv-1\Rightarrow n$ prim}) ist $9$ nicht prim.
\schritt{(c)}
Ist $n = ab$ mit $1 < a < b < n$, so stehen $a$ und $b$ beide in $(n - 1)!$, also $n\mid(n - 1)!$. Ist $n = p^{2}$ mit $p > 2$, so stehen $p$ und $2p < p^{2}$ darin, also $p^{2}\mid(n - 1)!$. Ausnahme $n = 4$: $3! = 6\equiv2\pmod4$.
\begin{mitnehmen}\kastentext
Wilson ist eine Äquivalenz --- im Gegensatz zu Fermat (Aufgabe 1.30) ein perfekter Primzahltest. Nur rechnet niemand $(n - 1)!$ für $n$ mit $300$ Stellen; die Fakultät hat keine schnelle Potenzierung.
\end{mitnehmen}
\end{loesungkarte}

%%%%%%%%%% 1.37
\stupserzwei{tief}{$\bigl(\frac3{41}\bigr) = \bigl(\frac{41}3\bigr) = \bigl(\frac23\bigr) = -1$ und $\bigl(\frac5{41}\bigr) = \bigl(\frac{41}5\bigr) = \bigl(\frac15\bigr) = 1$.}{22pt}
\begin{loesungkarte}{tief}{1.37}{Ein Legendre-Symbol, Schritt für Schritt}
\derweg{Multiplikativität im Zähler, Reziprozität zum Umdrehen, Reduktion des Zählers modulo Nenner, Ergänzungssätze für $-1$ und $2$.}
\schritt{(a)}
$\bigl(\frac{15}{41}\bigr) = \bigl(\frac3{41}\bigr)\bigl(\frac5{41}\bigr)$. $41\equiv1\pmod4$, also dreht Reziprozität ohne Vorzeichen: $\bigl(\frac3{41}\bigr) = \bigl(\frac{41}3\bigr) = \bigl(\frac23\bigr) = -1$ ($2$ ist kein Quadrat modulo $3$); $\bigl(\frac5{41}\bigr) = \bigl(\frac{41}5\bigr) = \bigl(\frac15\bigr) = 1$. Produkt: $\bigl(\frac{15}{41}\bigr) = -1$. $x^{2}\equiv15\pmod{41}$ ist unlösbar.
\schritt{(b)}
Euler: $15^{20}\bmod41$. $15^{2} = 225 = 5\cdot41 + 20\equiv20$; $15^{4}\equiv400 = 9\cdot41 + 31\equiv31\equiv-10$; $15^{8}\equiv100\equiv18$; $15^{16}\equiv324 = 7\cdot41 + 37\equiv37\equiv-4$; $15^{20} = 15^{16}\cdot15^{4}\equiv(-4)(-10) = 40\equiv-1$. \checkmark
\schritt{(c)}
$\bigl(\frac2{41}\bigr) = +1$, da $41\equiv1\pmod8$. $\bigl(\frac{-1}{41}\bigr) = +1$, da $41\equiv1\pmod4$. (Tatsächlich $9^{2} = 81\equiv-1$ und $17^{2} = 289 = 7\cdot41 + 2\equiv2$.)
\begin{mitnehmen}\kastentext
Legendre-Symbole rechnet man wie einen euklidischen Algorithmus: umdrehen, reduzieren, Zweierpotenzen abspalten. Jeder Schritt braucht nur \enquote{modulo $4$} und \enquote{modulo $8$} --- nie eine Quadratzahl.
\end{mitnehmen}
\end{loesungkarte}

%%%%%%%%%% 1.38
\stupserzwei{gruen}{Quadrate modulo $13$: $1, 4, 9, 3, 12, 10$. Die $5$ fehlt.}{22pt}
\begin{loesungkarte}{gruen}{1.38}{$x^{2}\equiv5\pmod{13}$}
\derweg{Einmal die sechs Quadrate hinschreiben, einmal Reziprozität.}
\schritt{(a)}
$x^{2}$ für $x = 1, \dots, 6$ (die anderen sind Negative): $1, 4, 9, 16\equiv3, 25\equiv12, 36\equiv10$. Quadratische Reste: $\{1, 3, 4, 9, 10, 12\}$. $5$ ist keiner --- unlösbar.
\schritt{(b)}
$\bigl(\frac5{13}\bigr) = \bigl(\frac{13}5\bigr)$ ($13\equiv1\pmod4$) $= \bigl(\frac35\bigr) = \bigl(\frac53\bigr)$ ($5\equiv1\pmod4$) $= \bigl(\frac23\bigr) = -1$. \checkmark
\begin{mitnehmen}\kastentext
Genau die Hälfte der Reste $1, \dots, p - 1$ sind Quadrate. Für $p = 13$ kann man sie hinschreiben; für $p = 383$ (nächste Aufgabe) nicht mehr --- und das Reziprozitätsgesetz merkt den Unterschied nicht.
\end{mitnehmen}
\end{loesungkarte}

%%%%%%%%%% 1.39
\stupserzwei{gruen}{$\bigl(\frac3{383}\bigr) = -\bigl(\frac{383}3\bigr) = -\bigl(\frac23\bigr) = +1$; $\bigl(\frac{73}{383}\bigr) = \bigl(\frac{383}{73}\bigr) = \bigl(\frac{18}{73}\bigr) = \bigl(\frac2{73}\bigr)\bigl(\frac9{73}\bigr)$.}{22pt}
\begin{loesungkarte}{gruen}{1.39}{$x^{2}\equiv219\pmod{383}$}
\derweg{Zähler faktorisieren, für jeden Faktor Reziprozität mit dem richtigen Vorzeichen, Ergänzungssatz für $2$.}
\schritt{Rechnung}
$219 = 3\cdot73$, $383\equiv3\pmod4$.
\begin{itemize}[leftmargin=1.2em,itemsep=1pt]
\item $\bigl(\frac3{383}\bigr)$: $3\equiv3$ und $383\equiv3\pmod4$, also Vorzeichenwechsel: $= -\bigl(\frac{383}3\bigr) = -\bigl(\frac23\bigr) = -(-1) = +1$.
\item $\bigl(\frac{73}{383}\bigr)$: $73\equiv1\pmod4$, kein Vorzeichen: $= \bigl(\frac{383}{73}\bigr)$; $383 = 5\cdot73 + 18$, also $= \bigl(\frac{18}{73}\bigr) = \bigl(\frac2{73}\bigr)\bigl(\frac{9}{73}\bigr) = \bigl(\frac2{73}\bigr)\cdot1$; $73\equiv1\pmod8$, also $\bigl(\frac2{73}\bigr) = +1$.
\end{itemize}
$\bigl(\frac{219}{383}\bigr) = (+1)(+1) = +1$: \textbf{lösbar}.
\probe{Tatsächlich $169^{2} = 28\,561 = 74\cdot383 + 219$. \checkmark\ (Die zweite Lösung ist $383 - 169 = 214$.)}
\begin{mitnehmen}\kastentext
Das Reziprozitätsgesetz entscheidet die Lösbarkeit, ohne eine Lösung zu nennen --- $169$ fand ein Rechner. Die Lösung selbst zu finden ist ein eigenes Problem (Tonelli--Shanks), das für $p\equiv3\pmod4$ einfach ist: $x\equiv a^{(p+1)/4}$.
\end{mitnehmen}
\end{loesungkarte}

%%%%%%%%%% 1.40
\stupserzwei{gruen}{$\bigl(\frac{1001}{9907}\bigr) = \bigl(\frac{9907}{1001}\bigr) = \bigl(\frac{898}{1001}\bigr) = \bigl(\frac2{1001}\bigr)\bigl(\frac{449}{1001}\bigr)$, und dann weiter umdrehen --- der Zähler wird jedes Mal kleiner.}{22pt}
\begin{loesungkarte}{gruen}{1.40}{1001 über 9907}
\derweg{Jacobi: umdrehen (Vorzeichen nur, wenn beide $\equiv3\pmod4$), modulo reduzieren, Zweien abspalten --- wie Euklid, ohne zu faktorisieren.}
\schritt{(a)}
$1001\equiv1\pmod4$, also $\bigl(\frac{1001}{9907}\bigr) = \bigl(\frac{9907}{1001}\bigr)$. $9907 = 9\cdot1001 + 898$: $= \bigl(\frac{898}{1001}\bigr) = \bigl(\frac2{1001}\bigr)\bigl(\frac{449}{1001}\bigr)$. $1001\equiv1\pmod8$: $\bigl(\frac2{1001}\bigr) = 1$. $449\equiv1\pmod4$: $\bigl(\frac{449}{1001}\bigr) = \bigl(\frac{1001}{449}\bigr) = \bigl(\frac{103}{449}\bigr)$ ($1001 - 2\cdot449 = 103$). $449\equiv1\pmod4$: $= \bigl(\frac{449}{103}\bigr) = \bigl(\frac{37}{103}\bigr)$ ($449 - 4\cdot103 = 37$). $37\equiv1\pmod4$: $= \bigl(\frac{103}{37}\bigr) = \bigl(\frac{29}{37}\bigr)$. $29\equiv1\pmod4$: $= \bigl(\frac{37}{29}\bigr) = \bigl(\frac8{29}\bigr) = \bigl(\frac2{29}\bigr)^{3} = \bigl(\frac2{29}\bigr)$; $29\equiv5\pmod8$: $= -1$.
\[ \Bigl(\frac{1001}{9907}\Bigr) = -1 . \]
\schritt{(b)}
$1001 = 7\cdot11\cdot13$; $9907\equiv3\pmod4$. $\bigl(\frac7{9907}\bigr) = -\bigl(\frac{9907}7\bigr) = -\bigl(\frac27\bigr) = -1$ ($7\equiv7\pmod8$). $\bigl(\frac{11}{9907}\bigr) = -\bigl(\frac{9907}{11}\bigr) = -\bigl(\frac7{11}\bigr) = -\bigl(-\bigl(\frac{11}7\bigr)\bigr) = \bigl(\frac47\bigr) = 1$. $\bigl(\frac{13}{9907}\bigr) = \bigl(\frac{9907}{13}\bigr) = \bigl(\frac1{13}\bigr) = 1$. Produkt $-1$. \checkmark\ Derselbe Wert --- aber man musste $1001$ faktorisieren, was bei großen Zahlen unmöglich ist; der Jacobi-Weg braucht nur Divisionen mit Rest.
\schritt{(c)}
$9907$ ist prim, also ist das Jacobi-Symbol ein Legendre-Symbol: $x^{2}\equiv1001\pmod{9907}$ ist \textbf{unlösbar}. Für zusammengesetztes $N$ dagegen heißt $\bigl(\frac aN\bigr) = 1$ \emph{nichts} über die Lösbarkeit: $\bigl(\frac2{15}\bigr) = \bigl(\frac23\bigr)\bigl(\frac25\bigr) = (-1)(-1) = 1$, aber $2$ ist kein Quadrat modulo $15$ (nicht einmal modulo $3$). Nur $\bigl(\frac aN\bigr) = -1$ ist aussagekräftig: dann ist $a$ sicher kein Quadrat.
\begin{mitnehmen}\kastentext
Das Jacobi-Symbol ist der Preis dafür, Reziprozität ohne Faktorisierung zu rechnen: Es gehorcht denselben Regeln, aber $+1$ beweist nichts mehr. Genau diese Eigenschaft nutzt der Solovay--Strassen-Primzahltest (Kapitel 18).
\end{mitnehmen}
\end{loesungkarte}

%%%%%%%%%% 1.41
\stupserzwei{tief}{$\sigma(28) = \sigma(4)\sigma(7) = 7\cdot8 = 56 = 2\cdot28$. Für (b): $\sigma(2^{p-1}) = 2^{p} - 1$ und $\sigma(2^{p} - 1) = 2^{p}$, wenn $2^{p} - 1$ prim ist.}{22pt}
\begin{loesungkarte}{tief}{1.41}{Vollkommen, multiplikativ}
\derweg{$\sigma$ über die Zerlegung; die Euklid-Formel folgt aus zwei Zeilen Multiplikativität.}
\schritt{(a)}
$28 = 2^{2}\cdot7$: $\sigma(28) = (1 + 2 + 4)(1 + 7) = 7\cdot8 = 56 = 2\cdot28$. $496 = 2^{4}\cdot31$: $\sigma = 31\cdot32 = 992 = 2\cdot496$. Beide vollkommen.
\schritt{(b)}
$N = 2^{p-1}(2^{p} - 1)$ mit $q = 2^{p} - 1$ prim; $2^{p-1}$ und $q$ sind teilerfremd ($q$ ungerade). $\sigma(N) = \sigma(2^{p-1})\sigma(q) = (2^{p} - 1)(q + 1) = (2^{p} - 1)\cdot2^{p} = 2N$. $\blacksquare$
\schritt{(c)}
Jeder Teiler von $ab$ ist eindeutig $d_1d_2$ mit $d_1\mid a$, $d_2\mid b$ (Fundamentalsatz: die Primfaktoren von $a$ und $b$ sind disjunkt). Also $\sigma(ab) = \sum_{d_1\mid a}\sum_{d_2\mid b}d_1d_2 = \bigl(\sum d_1\bigr)\bigl(\sum d_2\bigr) = \sigma(a)\sigma(b)$. $\blacksquare$
\schritt{(d)}
Nein: $\sigma(4) = 7$, aber $\sigma(2)^{2} = 9$. Die Teiler von $4$ sind $1, 2, 4$ --- der Teiler $2$ lässt sich auf zwei Arten als $d_1d_2$ schreiben ($1\cdot2$ und $2\cdot1$), die Zerlegung ist nicht mehr eindeutig. Multiplikativ heißt: nur für teilerfremde Faktoren.
\begin{mitnehmen}\kastentext
Multiplikative Funktionen ($\tau, \sigma, \varphi, \mu$) sind durch ihre Werte auf Primzahlpotenzen bestimmt. Die Euklid-Formel $2^{p-1}(2^{p} - 1)$ ist eine Zeile $\sigma$-Rechnung; dass sie \emph{alle} geraden vollkommenen Zahlen liefert, ist Eulers Satz und schwerer (Band 1, Kapitel 9.6).
\end{mitnehmen}
\end{loesungkarte}

%%%%%%%%%% 1.42
\stupserzwei{gruen}{$\sigma(2^{4}) = 31$, $\sigma(3^{2}) = 13$, $\sigma(5) = 6$.}{22pt}
\begin{loesungkarte}{gruen}{1.42}{$\tau(720)$ und $\sigma(720)$}
\derweg{Zerlegung, Formeln, multiplizieren.}
\schritt{Rechnung}
$720 = 2^{4}\cdot3^{2}\cdot5$. $\tau(720) = 5\cdot3\cdot2 = 30$. $\sigma(720) = 31\cdot13\cdot6 = 2418$. Da $2418 > 1440 = 2\cdot720$: abundant --- sogar sehr: $\sigma(720)/720 = 3{,}36$.
\begin{mitnehmen}\kastentext
Zahlen mit vielen kleinen Primfaktoren sind abundant; Primzahlen und ihre Potenzen defizient ($\sigma(p^{k}) < 2p^{k}$). Vollkommen ist die Grenze --- und ob es ungerade vollkommene Zahlen gibt, weiß niemand.
\end{mitnehmen}
\end{loesungkarte}

%%%%%%%%%% 1.43
\stupserzwei{gruen}{$W(n) = \sum_{d\mid n}\mu(d)k^{n/d}$. Für $n = 4$, $k = 2$: $2^{4} - 2^{2} = 12$ primitive Wörter, in $12/4 = 3$ Halsketten.}{22pt}
\begin{loesungkarte}{gruen}{1.43}{Primitive Halsketten}
\derweg{Alle Wörter nach ihrer kleinsten Periode sortieren; Möbius-Inversion; durch $n$ teilen, weil jede primitive Halskette aus genau $n$ verschiedenen Wörtern besteht.}
\schritt{(a)}
Jedes Wort $w$ der Länge $n$ hat eine kleinste Periode $d$; dann ist $w$ die $n/d$-fache Wiederholung eines primitiven Wortes der Länge $d$, und $d\mid n$ (sonst wäre $\operatorname{ggT}(d, n)$ eine kleinere Periode). Umgekehrt gibt jedes primitive Wort der Länge $d\mid n$ genau ein Wort der Länge $n$. Also $k^{n} = \sum_{d\mid n}W(d)$. $\blacksquare$
\schritt{(b)}
Möbius-Inversion: $W(n) = \sum_{d\mid n}\mu(d)\,k^{n/d}$. Ein primitives Wort hat $n$ verschiedene zyklische Drehungen (hätte es weniger, hätte es eine echte Periode), also zerfallen die $W(n)$ primitiven Wörter in Halsketten zu je $n$: $P(n) = W(n)/n$.
\schritt{(c)}
$k = 2$, $n = 4$: $W(4) = 2^{4} - 2^{2} = 12$ ($\mu(1) = 1$, $\mu(2) = -1$, $\mu(4) = 0$), $P(4) = 3$: die Halsketten $aaab$, $aabb$, $abbb$. ($aaaa$, $bbbb$, $abab$ sind nicht primitiv.) \checkmark\ $n = 6$: $W(6) = 2^{6} - 2^{3} - 2^{2} + 2 = 64 - 8 - 4 + 2 = 54$, $P(6) = 9$.
\schritt{(d)}
Für $p$ prim: $W(p) = k^{p} - k$, und $P(p) = (k^{p} - k)/p$ ist eine Anzahl, also ganz: $p\mid k^{p} - k$. Das ist der kleine Fermat --- bewiesen durch Zählen von Halsketten.
\begin{mitnehmen}\kastentext
Möbius-Inversion: Wer $g(n) = \sum_{d\mid n}f(d)$ kennt, kennt $f(n) = \sum_{d\mid n}\mu(d)g(n/d)$. Die Halsketten sind das Lehrbuchbeispiel, und der kombinatorische Fermat-Beweis ist ein Geschenk am Rande.
\end{mitnehmen}
\end{loesungkarte}

%%%%%%%%%% 1.44
\stupserzwei{gruen}{$p\equiv1\pmod4$: $2^{p-1} = 16^{k}\equiv6\pmod{10}$ und $2^{p} - 1\equiv1$, Produkt $\equiv6$. $p\equiv3\pmod4$: zeigen Sie $16^{k}(8\cdot16^{k} - 1)\equiv7\pmod{25}$.}{22pt}
\begin{loesungkarte}{gruen}{1.44}{Gerade vollkommene Zahlen enden auf 6 oder 28}
\derweg{Euler--Euklid-Form einsetzen, $p$ modulo $4$ unterscheiden, Reste modulo $10$ bzw.\ $100$ verfolgen.}
\schritt{Fall $p = 2$}
$N = 6$. Endet auf $6$.
\schritt{Fall $p\equiv1\pmod4$, $p = 4k + 1$}
$2^{p-1} = 2^{4k} = 16^{k}\equiv6^{k}\equiv6\pmod{10}$ (denn $6\cdot6 = 36\equiv6$). $2^{p} - 1 = 2\cdot16^{k} - 1\equiv12 - 1\equiv1\pmod{10}$. $N\equiv6\cdot1 = 6\pmod{10}$: endet auf $6$.
\schritt{Fall $p\equiv3\pmod4$, $p = 4k + 3$}
$2^{p-1} = 4\cdot16^{k}$, $2^{p} - 1 = 8\cdot16^{k} - 1$, also $N = 4\cdot16^{k}(8\cdot16^{k} - 1)$. Modulo $25$ hat $16^{k}$ die Periode $5$: $16^{k}\equiv1, 16, 6, 21, 11$ für $k\equiv0, 1, 2, 3, 4\pmod5$. Entsprechend $8\cdot16^{k} - 1\equiv7, 2, 22, 17, 12$. Die Produkte $16^{k}(8\cdot16^{k} - 1)$: $7, 32, 132, 357, 132$ --- alle $\equiv7\pmod{25}$. Also $N\equiv4\cdot7 = 28\pmod{100}$ (denn $4\cdot(7 + 25j) = 28 + 100j$): endet auf $28$. $\blacksquare$
\probe{$p = 3$: $28$. $p = 5$: $496$. $p = 7$: $8128$. $p = 13$: $33\,550\,336$. $p = 19$: $137\,438\,691\,328$. \checkmark}
\begin{mitnehmen}\kastentext
Die Endziffern sind kein Zufall, sondern Restrechnung mit einer Formel. Die Fallunterscheidung modulo $4$ kommt daher, dass $2^{4} = 16\equiv6\pmod{10}$ \enquote{stabil} ist --- $6^{k}\equiv6$.
\end{mitnehmen}
\end{loesungkarte}

%%%%%%%%%% 1.45
\stupserzwei{akzent}{$2047 = 23\cdot89$. Und $2^{ab} - 1 = (2^{a} - 1)(2^{a(b-1)} + \dots + 2^{a} + 1)$.}{22pt}
\begin{loesungkarte}{akzent}{1.45}{$2^{11} - 1$}
\derweg{Ein Gegenbeispiel, die richtige Richtung des Satzes, und der Grund, warum die Teiler so aussehen, wie sie aussehen.}
\schritt{(a)}
$2047 = 23\cdot89$. Zusammengesetzt --- die erste Mersenne-Zahl $M_p$ mit primem $p$, die nicht prim ist.
\schritt{(b)}
$x^{b} - 1 = (x - 1)(x^{b-1} + \dots + x + 1)$ mit $x = 2^{a}$: $2^{ab} - 1 = (2^{a} - 1)(2^{a(b-1)} + \dots + 1)$. Ist $n = ab$ mit $1 < a < n$, so ist $2^{a} - 1$ ein echter Teiler von $2^{n} - 1$. Also: $2^{n} - 1$ prim $\Rightarrow$ $n$ prim. Die Umkehrung ist (a). $\blacksquare$
\schritt{(c)}
Sei $q$ prim, $q\mid2^{p} - 1$, also $2^{p}\equiv1\pmod q$. Die Ordnung von $2$ modulo $q$ teilt $p$, ist also $1$ oder $p$; sie ist nicht $1$ ($2\not\equiv1$), also $p$. Nach Fermat teilt die Ordnung $q - 1$: $p\mid q - 1$. Da $q - 1$ gerade und $p$ ungerade ist, $2p\mid q - 1$: $q = 2kp + 1$. $\blacksquare$ Probe: $23 = 2\cdot11 + 1$, $89 = 8\cdot11 + 1$; $47 = 2\cdot23 + 1$, $178\,481 = 7760\cdot23 + 1$. \checkmark
\begin{fehler}\kastentext
\enquote{$2^{p} - 1$ ist prim für primes $p$.} --- Vier Beispiele ($p = 2, 3, 5, 7$), dann $2047 = 23\cdot89$. Richtig ist nur die Umkehrung. Bis heute kennt man $52$ Mersenne-Primzahlen; ob es unendlich viele sind, ist offen.
\end{fehler}
\begin{mitnehmen}\kastentext
Teiler von $2^{p} - 1$ sind $\equiv1\pmod{2p}$ --- das macht die Suche nach ihnen schnell und ist der Kern des Lucas--Lehmer-Tests, mit dem die größten bekannten Primzahlen gefunden werden (Band 1, Kapitel 9.6).
\end{mitnehmen}
\end{loesungkarte}

%%%%%%%%%% 1.46
\stupserzwei{gruen}{$\ln10^{20} = 20\ln10 = 46{,}05$. Also etwa eine Primzahl unter $46$ Zahlen.}{22pt}
\begin{loesungkarte}{gruen}{1.46}{Primzahlen jenseits von $10^{20}$}
\derweg{Primzahlsatz als Dichte $1/\ln x$, mal Intervalllänge.}
\schritt{(a)}
Dichte $1/\ln(10^{20}) = 1/46{,}05$. Erwartet $10^{6}/46{,}05\approx21\,700$ Primzahlen. (Das Intervall ist so kurz, dass $\ln x$ darin konstant ist.)
\schritt{(b)}
Im Mittel $46$ Zahlen pro Primzahl. Nur ungerade: $23$. Ohne Vielfache von $2, 3, 5$ bleiben $\frac12\cdot\frac23\cdot\frac45 = \frac4{15}$ aller Zahlen, also etwa $46\cdot\frac4{15}\approx12$ Kandidaten pro Primzahl.
\schritt{(c)}
$\ln10^{300} = 690{,}8$: etwa $691$ zufällige Zahlen, $345$ ungerade, $184$ nach Aussieben von $2, 3, 5$. Jeder Kandidat braucht einen Miller--Rabin-Test (Aufgabe 1.58) --- Millisekunden. RSA-Schlüssel zu erzeugen ist deshalb leicht; sie zu brechen nicht.
\begin{mitnehmen}\kastentext
Primzahlen sind häufig genug: Bei $n$ Stellen findet man eine unter etwa $2{,}3n$ Zahlen. Darauf beruht, dass Kryptographie mit Primzahlen praktikabel ist.
\end{mitnehmen}
\end{loesungkarte}

%%%%%%%%%% 1.47
\stupserzwei{gruen}{$p$, $p + 2$, $p + 4$ lassen modulo $3$ die Reste $r$, $r + 2$, $r + 1$ --- alle drei Reste kommen vor, einer ist $0$.}{22pt}
\begin{loesungkarte}{gruen}{1.47}{Der einzige Drilling}
\derweg{Restklassen modulo $3$ durchgehen.}
\schritt{(a)}
Von $p$, $p + 2$, $p + 4$ ist eine durch $3$ teilbar: Ist $p\equiv0$, dann $p$; ist $p\equiv1$, dann $p + 2$; ist $p\equiv2$, dann $p + 4$. Eine durch $3$ teilbare Primzahl ist $3$ selbst. Also ist eine der drei gleich $3$; da $p$ die kleinste ist, $p = 3$: $(3, 5, 7)$. $\blacksquare$
\schritt{(b)}
Zwillinge $(p, p + 2)$: Reste $r, r + 2$ --- der Rest $r + 1$ fehlt, also ist für $p\equiv2\pmod3$ keine der beiden durch $3$ teilbar ($5, 7$; $11, 13$; $17, 19$). Drillinge $(p, p + 2, p + 6)$: Reste $r, r + 2, r$ --- auch hier bleibt eine Restklasse frei ($5, 7, 11$; $11, 13, 17$; $17, 19, 23$). Solche \enquote{zulässigen} Muster vermeiden, dass eine Restklasse modulo jeder Primzahl vollständig besetzt ist, und von ihnen erwartet man unendlich viele Vorkommen (Hardy--Littlewood); bewiesen ist das für keines (Band 1, Kapitel 10.6).
\begin{mitnehmen}\kastentext
Ein Muster von Primzahlen kann nur dann oft vorkommen, wenn es modulo jeder Primzahl $q$ eine Restklasse frei lässt. $(p, p + 2, p + 4)$ belegt modulo $3$ alle drei --- deshalb nur einmal.
\end{mitnehmen}
\end{loesungkarte}

%%%%%%%%%% 1.48
\stupserzwei{gruen}{$m = 2\colon n = 1$. $m = 3\colon n = 2$. $m = 4\colon n = 1, 3$. $m = 5\colon n = 2, 4$. $m = 6\colon n = 1, 5$. $m = 7\colon n = 2, 4, 6$. $m = 8\colon n = 1, 3, 5$. $m = 9\colon n = 2, 4$ ($n = 8$ gibt $c = 145$).}{22pt}
\begin{loesungkarte}{gruen}{1.48}{Pythagoreische Tripel bis 100}
\derweg{$(m^{2} - n^{2}, 2mn, m^{2} + n^{2})$ für alle zulässigen Paare mit $m^{2} + n^{2}\le100$.}
\schritt{(a)}
\begin{center}
{\small\renewcommand{\arraystretch}{1.25}
\begin{tabular}{@{}cc|ccc@{}}
$m$ & $n$ & $a$ & $b$ & $c$\\\hline
2 & 1 & 3 & 4 & 5\\ 3 & 2 & 5 & 12 & 13\\ 4 & 1 & 15 & 8 & 17\\ 4 & 3 & 7 & 24 & 25\\ 5 & 2 & 21 & 20 & 29\\ 5 & 4 & 9 & 40 & 41\\ 6 & 1 & 35 & 12 & 37\\ 6 & 5 & 11 & 60 & 61\\
\end{tabular}\qquad
\begin{tabular}{@{}cc|ccc@{}}
$m$ & $n$ & $a$ & $b$ & $c$\\\hline
7 & 2 & 45 & 28 & 53\\ 7 & 4 & 33 & 56 & 65\\ 7 & 6 & 13 & 84 & 85\\ 8 & 1 & 63 & 16 & 65\\ 8 & 3 & 55 & 48 & 73\\ 8 & 5 & 39 & 80 & 89\\ 9 & 2 & 77 & 36 & 85\\ 9 & 4 & 65 & 72 & 97\\
\end{tabular}}
\end{center}
Sechzehn primitive Tripel mit $c\le100$.
\schritt{(b)}
Sind $m, n$ beide ungerade, so sind $m^{2} - n^{2}$, $2mn$, $m^{2} + n^{2}$ alle gerade --- das Tripel ist nicht primitiv (es ist das Doppelte eines anderen: $m = 3$, $n = 1$ gibt $(8, 6, 10) = 2\cdot(4, 3, 5)$). Haben $m, n$ einen gemeinsamen Primteiler, teilt er alle drei Zahlen. Beide Bedingungen zusammen sind genau die Primitivität (Band 1, Kapitel 11).
\schritt{(c)}
$65 = 5\cdot13$, beide $\equiv1\pmod4$, also $65 = 1^{2} + 8^{2} = 4^{2} + 7^{2}$ auf zwei Arten --- und jede Darstellung $c = m^{2} + n^{2}$ liefert ein Tripel: $(8, 1)\to(63, 16, 65)$, $(7, 4)\to(33, 56, 65)$. Ebenso $85 = 5\cdot17 = 2^{2} + 9^{2} = 6^{2} + 7^{2}$.
\begin{mitnehmen}\kastentext
Primitive Tripel entsprechen den Darstellungen der Hypotenuse als Summe zweier Quadrate --- ein Zusammenhang zwischen Kapitel 11 und 14, und der Grund, warum Hypotenusen immer Produkte von Primzahlen $\equiv1\pmod4$ sind (bis auf einen Faktor $2$ bei nicht-primitiven).
\end{mitnehmen}
\end{loesungkarte}

%%%%%%%%%% 1.49
\stupserzwei{tief}{$[3; 7, 15, 1, 292, \dots]$. Näherungsbrüche: $3$, $\frac{22}7$, $\frac{333}{106}$, $\frac{355}{113}$, $\frac{103\,993}{33\,102}$.}{22pt}
\begin{loesungkarte}{tief}{1.49}{$\pi$ als Kettenbruch}
\derweg{Ganzzahliger Anteil, Kehrwert des Rests, wiederholen; Näherungsbrüche mit der Rekursion $p_k = a_kp_{k-1} + p_{k-2}$.}
\schritt{(a)}
$\pi = 3{,}14159265$: $a_0 = 3$, Rest $0{,}14159265$, Kehrwert $7{,}0625133$: $a_1 = 7$, Rest $0{,}0625133$, Kehrwert $15{,}996594$: $a_2 = 15$, Rest $0{,}996594$, Kehrwert $1{,}0034172$: $a_3 = 1$, Rest $0{,}0034172$, Kehrwert $292{,}63$: $a_4 = 292$. Also $\pi = [3; 7, 15, 1, 292, 1, 1, 1, 2, \dots]$ --- ohne erkennbares Muster, im Gegensatz zu $\sqrt7$ (nächste Aufgabe) oder $e = [2; 1, 2, 1, 1, 4, 1, 1, 6, \dots]$.
\schritt{(b)}
\begin{center}
{\small\renewcommand{\arraystretch}{1.3}
\begin{tabular}{@{}lllll@{}}
$a_k$ & $3$ & $7$ & $15$ & $1$\\
$p_k/q_k$ & $3/1$ & $22/7$ & $333/106$ & $355/113$\\
Fehler & $0{,}14$ & $1{,}3\cdot10^{-3}$ & $8{,}3\cdot10^{-5}$ & $2{,}7\cdot10^{-7}$\\
\end{tabular}}
\end{center}
($p_2 = 15\cdot22 + 3 = 333$, $q_2 = 15\cdot7 + 1 = 106$; $p_3 = 333 + 22 = 355$, $q_3 = 106 + 7 = 113$.)
\schritt{(c)}
Der Fehler eines Näherungsbruchs ist etwa $\frac1{a_{k+1}q_k^{2}}$: Ein großer \emph{nächster} Teilnenner macht den \emph{aktuellen} Bruch außergewöhnlich gut. Nach $355/113$ kommt $a_4 = 292$, also Fehler $\approx\frac1{292\cdot113^{2}} = 2{,}7\cdot10^{-7}$ --- sechs richtige Dezimalen mit einem dreistelligen Nenner. Der nächste Näherungsbruch $103\,993/33\,102$ ist kaum besser und viel unhandlicher. Zu Chongzhi fand $355/113$ im 5.\ Jahrhundert; in Europa Adriaan Anthonisz um 1585.
\begin{mitnehmen}\kastentext
Näherungsbrüche sind die besten Näherungen überhaupt: Kein Bruch mit kleinerem Nenner kommt näher. Und ein großer Teilnenner ist ein Geschenk --- er sagt, dass der Bruch davor \enquote{zu gut für seinen Nenner} ist.
\end{mitnehmen}
\end{loesungkarte}

%%%%%%%%%% 1.50
\stupserzwei{gruen}{$\sqrt7 = [2; \overline{1, 1, 1, 4}]$. Die Periode $4$ endet mit $2a_0 = 4$ --- das ist immer so.}{22pt}
\begin{loesungkarte}{gruen}{1.50}{Der Kettenbruch von $\sqrt7$}
\derweg{Der Algorithmus mit $\sqrt7$ exakt: Rest, Kehrwert rational machen, Ganzteil, bis sich ein Rest wiederholt.}
\schritt{(a)}
$\sqrt7 = 2 + (\sqrt7 - 2)$. $\frac1{\sqrt7 - 2} = \frac{\sqrt7 + 2}{3} = 1 + \frac{\sqrt7 - 1}3$. $\frac3{\sqrt7 - 1} = \frac{3(\sqrt7 + 1)}6 = \frac{\sqrt7 + 1}2 = 1 + \frac{\sqrt7 - 1}2$. $\frac2{\sqrt7 - 1} = \frac{2(\sqrt7 + 1)}6 = \frac{\sqrt7 + 1}3 = 1 + \frac{\sqrt7 - 2}3$. $\frac3{\sqrt7 - 2} = \frac{3(\sqrt7 + 2)}3 = \sqrt7 + 2 = 4 + (\sqrt7 - 2)$. Der Rest $\sqrt7 - 2$ ist wieder da: $\sqrt7 = [2; \overline{1, 1, 1, 4}]$.
\schritt{(b)}
$a_k$: $2, 1, 1, 1, 4, 1, \dots$ Näherungsbrüche: $\frac21, \frac31, \frac52, \frac83, \frac{37}{14}, \frac{45}{17}, \frac{82}{31}, \dots$ ($\frac{37}{14} = 2{,}642857$ gegen $\sqrt7 = 2{,}645751$.)
\schritt{(c)}
$p^{2} - 7q^{2}$: $4 - 7 = -3$; $9 - 7 = 2$; $25 - 28 = -3$; $64 - 63 = 1$; $1369 - 1372 = -3$; $2025 - 2023 = 2$; $6724 - 6727 = -3$; dann wieder $1$ bei $\frac{127}{48}$. Die Werte sind klein ($|p^{2} - 7q^{2}| < 2\sqrt7$) und periodisch wie der Kettenbruch --- und am Ende jeder Periode steht $1$.
\begin{mitnehmen}\kastentext
Lagrange: Der Kettenbruch einer quadratischen Irrationalzahl ist periodisch. Die Näherungsbrüche am Periodenende lösen $p^{2} - dq^{2} = \pm1$ --- die Pellsche Gleichung fällt aus dem Kettenbruch heraus.
\end{mitnehmen}
\end{loesungkarte}

%%%%%%%%%% 1.51
\stupserzwei{gruen}{$(8 + 3\sqrt7)^{2} = 64 + 48\sqrt7 + 63 = 127 + 48\sqrt7$.}{22pt}
\begin{loesungkarte}{gruen}{1.51}{$x^{2} - 7y^{2} = 1$}
\derweg{Grundlösung aus dem Kettenbruch, weitere durch Potenzieren von $x + y\sqrt7$.}
\schritt{(a)}
Aus Aufgabe 1.50c: $(x, y) = (8, 3)$, denn $64 - 63 = 1$. Kleinere positive Lösungen gibt es nicht (die Näherungsbrüche davor geben $-3$ und $2$, und jede Lösung ist ein Näherungsbruch).
\schritt{(b)}
$(8 + 3\sqrt7)^{2} = 64 + 2\cdot24\sqrt7 + 9\cdot7 = 127 + 48\sqrt7$: $(127, 48)$, Probe $16129 - 7\cdot2304 = 16129 - 16128 = 1$. \checkmark
\schritt{(c)}
Für $\alpha = x + y\sqrt7$ sei $\bar\alpha = x - y\sqrt7$ und $N(\alpha) = \alpha\bar\alpha = x^{2} - 7y^{2}$. Die Norm ist multiplikativ: $N(\alpha\beta) = \alpha\beta\bar\alpha\bar\beta = N(\alpha)N(\beta)$. Also $N(\alpha^{n}) = N(\alpha)^{n} = 1$, und $\alpha^{n} = x_n + y_n\sqrt7$ liefert eine Lösung $(x_n, y_n)$. Umgekehrt ist jede Lösung eine Potenz der Grundlösung (Band 1, Kapitel 13.2).
\begin{mitnehmen}\kastentext
Die Lösungen der Pellschen Gleichung bilden eine Gruppe --- die Einheiten von $\mathbb Z[\sqrt7]$ mit Norm $1$ ---, erzeugt von einer Grundlösung. Für $d = 61$ ist die Grundlösung $x = 1\,766\,319\,049$: Fermat stellte sie 1657 als Herausforderung, weil er wusste, wie groß sie ist.
\end{mitnehmen}
\end{loesungkarte}

%%%%%%%%%% 1.52
\stupserzwei{akzent}{Quadrate bis $21$: $0, 1, 4, 9, 16$; keine zwei summieren sich zu $21$. Und $21 = 3\cdot7$ mit $3, 7\equiv3\pmod4$ in ungerader Potenz.}{22pt}
\begin{loesungkarte}{akzent}{1.52}{21 ist $\equiv1\pmod4$}
\derweg{Nachrechnen, dann den Satz richtig formulieren: Er handelt von den Exponenten der Primteiler $\equiv3\pmod4$.}
\schritt{(a)}
$21 - 0 = 21$, $21 - 1 = 20$, $21 - 4 = 17$, $21 - 9 = 12$, $21 - 16 = 5$: kein Quadrat. $21$ ist keine Summe zweier Quadrate --- obwohl $21\equiv1\pmod4$.
\schritt{(b)}
$45 = 36 + 9 = 6^{2} + 3^{2}$. Der Teiler $3\equiv3\pmod4$ kommt mit \emph{geradem} Exponenten vor: $3^{2}\cdot5$, und $5 = 1 + 4$; $45 = 3^{2}(1^{2} + 2^{2}) = 3^{2} + 6^{2}$.
\schritt{(c)}
$n$ ist Summe zweier Quadrate genau dann, wenn in der Primfaktorzerlegung jede Primzahl $\equiv3\pmod4$ mit geradem Exponenten vorkommt. (Primzahlen $\equiv1\pmod4$ und die $2$ sind unbeschränkt erlaubt.) Fermats Satz \enquote{$p\equiv1\pmod4\Rightarrow p = x^{2} + y^{2}$} ist der Baustein für Primzahlen; die Identität $(a^{2} + b^{2})(c^{2} + d^{2}) = (ac - bd)^{2} + (ad + bc)^{2}$ setzt zusammen.
\schritt{(d)}
Quadrate sind $\equiv0$ oder $1\pmod4$; Summen zweier Quadrate also $\equiv0, 1, 2$ --- nie $3$. Das zeigt sofort $3, 7, 11, 19\notin$. Aber $21\equiv1$ wird davon nicht erfasst: Die Bedingung modulo $4$ ist notwendig, nicht hinreichend.
\begin{fehler}\kastentext
\enquote{$n\equiv1\pmod4$, also Summe zweier Quadrate.} --- Der Satz von Fermat gilt für Primzahlen. Für zusammengesetzte Zahlen ist $21 = 3\cdot7$ das Gegenbeispiel: zwei schlechte Primteiler, beide einmal, das Produkt $\equiv1$, und trotzdem nicht darstellbar.
\end{fehler}
\begin{mitnehmen}\kastentext
Sätze über Primzahlen übertragen sich auf zusammengesetzte Zahlen nur über die Zerlegung --- und dann zählen Exponenten. Dasselbe Muster bei $\varphi$, bei quadratischen Resten (Jacobi, Aufgabe 1.40c) und beim Fermat-Test (Aufgabe 1.30).
\end{mitnehmen}
\end{loesungkarte}

%%%%%%%%%% 1.53
\stupserzwei{gruen}{$13 = 9 + 4$, $17 = 16 + 1$; Brahmagupta: $(ac\mp bd)^{2} + (ad\pm bc)^{2}$ mit zwei Vorzeichenwahlen gibt zwei Darstellungen von $221$.}{22pt}
\begin{loesungkarte}{gruen}{1.53}{Summen zweier Quadrate}
\derweg{Zerlegen, Primteiler $\equiv3\pmod4$ auf gerade Exponenten prüfen, Darstellungen zusammensetzen.}
\schritt{221}
$221 = 13\cdot17$, beide $\equiv1\pmod4$: ja. $13 = 3^{2} + 2^{2}$, $17 = 4^{2} + 1^{2}$. Brahmagupta mit $(a, b, c, d) = (3, 2, 4, 1)$: $(12 - 2)^{2} + (3 + 8)^{2} = 10^{2} + 11^{2}$ und $(12 + 2)^{2} + (3 - 8)^{2} = 14^{2} + 5^{2}$. Also $221 = 100 + 121 = 196 + 25$ --- zwei Darstellungen, weil zwei verschiedene Primteiler $\equiv1\pmod4$.
\schritt{198}
$198 = 2\cdot3^{2}\cdot11$. $11\equiv3\pmod4$ mit Exponent $1$: nein. (Probe: $198 - k^{2}$ für $k\le14$: $198, 197, 194, 189, 182, 173, 162, 149, 134, 117, 98, 77, 54, 29, 2$ --- kein Quadrat. \checkmark)
\schritt{2025}
$2025 = 3^{4}\cdot5^{2}$. $3$ mit geradem Exponenten: ja. Darstellungen: $45^{2} + 0^{2}$ und, aus $25 = 3^{2} + 4^{2}$ mal $9^{2}$: $27^{2} + 36^{2} = 729 + 1296$. \checkmark\ Das sind alle (mit $x\le y$).
\begin{mitnehmen}\kastentext
Entscheiden über die Zerlegung, darstellen über Brahmaguptas Identität --- die Zahl der Darstellungen wächst mit der Zahl der Primteiler $\equiv1\pmod4$ (Jacobis Formel zählt sie exakt).
\end{mitnehmen}
\end{loesungkarte}

%%%%%%%%%% 1.54
\stupserzwei{gruen}{$\lambda = \frac{3\cdot4 - 2}{2\cdot3}$? Nein --- hier ist $a = 0$: $\lambda = \frac{3\cdot(-2)^{2}}{2\cdot3} = 2$. Dann $x_3 = 4 + 4 = 8$, $y_3 = 2(-2 - 8) - 3 = -23$.}{22pt}
\begin{loesungkarte}{gruen}{1.54}{Ein Punkt wird verdoppelt}
\derweg{Tangentenformel für $2P$, Sekantenformel für $2P + P$; jeden Punkt an der Kurvengleichung prüfen.}
\schritt{(a)}
$y^{2} = x^{3} + 17$, also $a = 0$. $\lambda = \frac{3x_1^{2}}{2y_1} = \frac{12}6 = 2$. $x_3 = \lambda^{2} - 2x_1 = 4 + 4 = 8$, $y_3 = \lambda(x_1 - x_3) - y_1 = 2(-2 - 8) - 3 = -23$. $2P = (8, -23)$. Probe: $8^{3} + 17 = 529 = 23^{2}$. \checkmark
\schritt{(b)}
$2P = (8, -23)$, $P = (-2, 3)$: $\lambda = \frac{-23 - 3}{8 + 2} = -\frac{13}5$. $x_3 = \lambda^{2} - 8 - (-2) = \frac{169}{25} - 6 = \frac{19}{25}$. $y_3 = \lambda(x_1 - x_3) - y_1$ mit $(x_1, y_1) = (-2, 3)$: $-\frac{13}5\bigl(-2 - \frac{19}{25}\bigr) - 3 = -\frac{13}5\cdot\bigl(-\frac{69}{25}\bigr) - 3 = \frac{897}{125} - \frac{375}{125} = \frac{522}{125}$. $3P = \bigl(\frac{19}{25}, \frac{522}{125}\bigr)$. Probe: $\bigl(\frac{19}{25}\bigr)^{3} + 17 = \frac{6859 + 265\,625}{15\,625} = \frac{272\,484}{15\,625} = \bigl(\frac{522}{125}\bigr)^{2}$. \checkmark
\schritt{(c)}
Die Nenner wachsen ($1, 1, 25, \dots$) --- bei einem Punkt endlicher Ordnung müssten die Vielfachen irgendwann zum Ausgangspunkt zurückkehren, und nach Nagell--Lutz hätten Torsionspunkte ganzzahlige Koordinaten. $3P$ ist nicht ganzzahlig, also hat $P$ unendliche Ordnung: $E(\mathbb Q)$ ist unendlich, die Kurve hat positiven Rang.
\begin{mitnehmen}\kastentext
Das Gruppengesetz ist Rechnen mit Geraden und der Kurvengleichung --- jede Zeile prüfbar. Die Höhe (Größe der Nenner) wächst bei Punkten unendlicher Ordnung quadratisch; daran erkennt man den Rang. Band 29 hat die Geometrie dahinter.
\end{mitnehmen}
\end{loesungkarte}

%%%%%%%%%% 1.55
\stupserzwei{tief}{$3^{2} = 9$, $3^{4} = 81$, $3^{8} = 561$, $3^{16} = 721$, $3^{32} = 841$, $3^{64} = 281$, $3^{128} = 961$ (alles modulo $1000$). Bits von $218 = 11011010_2$: $128, 64, 16, 8, 2$.}{22pt}
\begin{loesungkarte}{tief}{1.55}{Quadrieren und Multiplizieren}
\derweg{Potenzen $3^{2^{j}}$ durch wiederholtes Quadrieren; die zum Exponenten gehörenden zusammenmultiplizieren --- alles modulo $1000$.}
\schritt{(a)}
Quadrate modulo $1000$: $3^{1} = 3$, $3^{2} = 9$, $3^{4} = 81$, $3^{8} = 6561\equiv561$, $3^{16}\equiv561^{2} = 314\,721\equiv721$, $3^{32}\equiv721^{2} = 519\,841\equiv841$, $3^{64}\equiv841^{2} = 707\,281\equiv281$, $3^{128}\equiv281^{2} = 78\,961\equiv961$.
$218 = 128 + 64 + 16 + 8 + 2$: $3^{218}\equiv961\cdot281\cdot721\cdot561\cdot9$. Schrittweise: $961\cdot281 = 270\,041\equiv41$; $41\cdot721 = 29\,561\equiv561$; $561\cdot561\equiv721$ (von oben); $721\cdot9 = 6489\equiv489$. Also $3^{218}\equiv489\pmod{1000}$.
\schritt{(b)}
Sieben Quadrierungen und vier Multiplikationen: elf. Naiv wären es $217$.
\schritt{(c)}
Höchstens $2047$ Quadrierungen und $2047$ Multiplikationen, also rund $4000$ --- gegen $2^{2048}$ naiv, eine Zahl mit $617$ Stellen. Deshalb ist $m^{e}\bmod n$ für RSA in Millisekunden zu berechnen.
\probe{$3^{218}\equiv489$: Kontrolle über $3^{200}\cdot3^{18}$; oder mit der Ordnung von $3$ modulo $1000$ ($100$): $3^{218}\equiv3^{18}$, und $3^{18} = 3^{16}\cdot3^{2}\equiv721\cdot9 = 6489\equiv489$. \checkmark}
\begin{mitnehmen}\kastentext
Schnelle Potenzierung: Binärdarstellung des Exponenten, so viele Quadrierungen wie Stellen, so viele Multiplikationen wie Einsen. Alle Zwischenergebnisse bleiben kleiner als $n^{2}$. Das ist der Algorithmus, den jede Chipkarte ausführt.
\end{mitnehmen}
\end{loesungkarte}

%%%%%%%%%% 1.56
\stupserzwei{gruen}{$A = 5^{6}\equiv8$, $B = 5^{15}\equiv19$. Schlüssel $= B^{6} = A^{15} = 2$.}{22pt}
\begin{loesungkarte}{gruen}{1.56}{Diffie--Hellman mit $p = 23$}
\derweg{Zwei Potenzen öffentlich, die gemeinsame Potenz privat; der Lauscher braucht den diskreten Logarithmus.}
\schritt{(a)}
$5^{2} = 25\equiv2$, also $5^{6} = 2^{3} = 8$: $A = 8$. $5^{15} = 5^{14}\cdot5 = 2^{7}\cdot5 = 128\cdot5 = 640 = 27\cdot23 + 19$: $B = 19$.
\schritt{(b)}
Alice: $B^{a} = 19^{6}$. $19\equiv-4$, $(-4)^{6} = 4096 = 178\cdot23 + 2$: $2$. Bob: $A^{b} = 8^{15} = 2^{45} = (2^{11})^{4}\cdot2$; $2^{11} = 2048 = 89\cdot23 + 1\equiv1$; also $2^{45}\equiv2$. Beide haben $2$. \checkmark
\schritt{(c)}
Der Lauscher muss $a$ aus $5^{a}\equiv8$ bestimmen (diskreter Logarithmus). Potenzen von $5$: $5, 2, 10, 4, 20, 8$ --- also $a = 6$ nach sechs Versuchen. Bei $p$ mit $600$ Stellen ist dieses Durchprobieren hoffnungslos, und bessere Verfahren (Index-Calculus) brauchen immer noch subexponentiell lange; genau diese Lücke zwischen \enquote{potenzieren ist leicht} und \enquote{Logarithmus ist schwer} ist das Verfahren.
\begin{mitnehmen}\kastentext
Diffie--Hellman (1976): $g^{ab}$ aus $g^{a}$ und $g^{b}$ berechnen kann nur, wer $a$ oder $b$ kennt --- solange der diskrete Logarithmus schwer ist. Heute meist auf elliptischen Kurven statt in $(\mathbb Z/p)^{\times}$, mit viel kleineren Zahlen.
\end{mitnehmen}
\end{loesungkarte}

%%%%%%%%%% 1.57
\stupserzwei{gruen}{$n = 3233$, $\varphi = 3120$, $d = 2753$ (denn $17\cdot2753 = 46\,801 = 15\cdot3120 + 1$), $c = 2790$.}{22pt}
\begin{loesungkarte}{gruen}{1.57}{RSA, vollständig}
\derweg{Schlüssel aus $p, q, e$; Potenzen mit Quadrieren und Multiplizieren; am Ende Euler.}
\schritt{(a)}
$n = 61\cdot53 = 3233$, $\varphi(n) = 60\cdot52 = 3120$. $d = 17^{-1}\bmod3120$: Euklid $3120 = 183\cdot17 + 9$, $17 = 1\cdot9 + 8$, $9 = 1\cdot8 + 1$; rückwärts $1 = 9 - 8 = 9 - (17 - 9) = 2\cdot9 - 17 = 2(3120 - 183\cdot17) - 17 = 2\cdot3120 - 367\cdot17$. Also $d\equiv-367\equiv2753$. Probe: $17\cdot2753 = 46\,801 = 15\cdot3120 + 1$. \checkmark
\schritt{(b)}
$c = 65^{17}\bmod3233$. $65^{2} = 4225\equiv992$; $65^{4}\equiv992^{2} = 984\,064\equiv984\,064 - 304\cdot3233 = 984\,064 - 982\,832 = 1232$; $65^{8}\equiv1232^{2} = 1\,517\,824\equiv1\,517\,824 - 469\cdot3233 = 1\,517\,824 - 1\,516\,277 = 1547$; $65^{16}\equiv1547^{2} = 2\,393\,209\equiv2\,393\,209 - 740\cdot3233 = 2\,393\,209 - 2\,392\,420 = 789$. $c = 65^{16}\cdot65\equiv789\cdot65 = 51\,285 = 15\cdot3233 + 2790$: $c = 2790$.
\schritt{(c)}
$2790^{2753}\bmod3233 = 65$ --- von Hand mit $12$ Quadrierungen und Multiplikationen (ein Rechner bestätigt $65$). Schneller: Mit $m^{d}\bmod61$ und $\bmod53$ und dem Restsatz (so machen es Chipkarten): $d\bmod60 = 53$, $2790\bmod61 = 45$, $45^{53}\bmod61 = 4$; $d\bmod52 = 49$, $2790\bmod53 = 34$, $34^{49}\bmod53 = 12$; und $x\equiv4\pmod{61}$, $x\equiv12\pmod{53}$ gibt $x = 65$. \checkmark
\schritt{(d)}
$ed\equiv1\pmod{\varphi(n)}$, also $ed = 1 + k\varphi(n)$, und $m^{ed} = m\cdot(m^{\varphi(n)})^{k}\equiv m\pmod n$ nach Euler, falls $\operatorname{ggT}(m, n) = 1$. (Für $\operatorname{ggT}(m, n) > 1$ hilft der Restsatz: modulo $p$ und $q$ getrennt gilt die Gleichung mit Fermat.) Wer $\varphi(n)$ kennt, kennt $d$; wer $\varphi(n)$ kennen will, muss $n$ faktorisieren.
\begin{mitnehmen}\kastentext
RSA ist der Satz von Euler, rückwärts gelesen: Der öffentliche Exponent potenziert, der private potenziert zurück, und zwischen beiden steht $\varphi(n)$, das nur kennt, wer $p$ und $q$ kennt.
\end{mitnehmen}
\end{loesungkarte}

%%%%%%%%%% 1.58
\stupserzwei{gruen}{$2^{35}\equiv263$, $263^{2}\equiv166$, $166^{2}\equiv67$, $67^{2}\equiv1$. Die $1$ kommt aus $67$, nicht aus $\pm1$ --- das ist der Beweis.}{22pt}
\begin{loesungkarte}{gruen}{1.58}{561 und Miller--Rabin}
\derweg{$n - 1 = 2^{s}t$ mit ungeradem $t$; $a^{t}$ berechnen, dann quadrieren und beobachten, ob die $1$ aus $-1$ entsteht.}
\schritt{(a)}
$560 = 2^{4}\cdot35$. $2^{35}\bmod561 = 263$ (Quadrieren und Multiplizieren). Quadrieren: $263^{2} = 69\,169 = 123\cdot561 + 166$, also $166$; $166^{2} = 27\,556 = 49\cdot561 + 67$, also $67$; $67^{2} = 4489 = 8\cdot561 + 1$, also $1$. Die Folge $263, 166, 67, 1$ erreicht $1$, aber der Wert davor ist $67\not\equiv\pm1$. Miller--Rabin: \textbf{zusammengesetzt}.
\schritt{(b)}
$67^{2}\equiv1$ mit $67\not\equiv\pm1$: eine nichttriviale Quadratwurzel von $1$ (Aufgabe 1.27). $561\mid(67 - 1)(67 + 1) = 66\cdot68$, also $\operatorname{ggT}(66, 561) = 33$ und $\operatorname{ggT}(68, 561) = 17$: $561 = 33\cdot17 = 3\cdot11\cdot17$. Der Test hat nicht nur entschieden, sondern faktorisiert.
\schritt{(c)}
$561 = 3\cdot11\cdot17$ und $560$ ist durch $2$, $10$ und $16$ teilbar, also durch $p - 1$ für jeden Primteiler: Nach Fermat ist $a^{560}\equiv1$ modulo $3$, $11$ und $17$, also modulo $561$, für jedes teilerfremde $a$ (Korselts Kriterium). Carmichael-Zahlen sind genau die mit dieser Eigenschaft; es gibt unendlich viele (Alford--Granville--Pomerance 1994).
\begin{mitnehmen}\kastentext
Fermat sieht nur $a^{n-1}$; Miller--Rabin sieht die Quadratwurzeln auf dem Weg dorthin. Eine Carmichael-Zahl täuscht den ersten Test für jede Basis, den zweiten für höchstens ein Viertel der Basen --- nach $20$ Basen ist die Irrtumswahrscheinlichkeit unter $10^{-12}$.
\end{mitnehmen}
\end{loesungkarte}

%%%%%%%%%% 1.59
\stupserzwei{gruen}{$90^{2} = 8100$, $8100 - 8051 = 49 = 7^{2}$. Also $8051 = 90^{2} - 7^{2} = 83\cdot97$.}{22pt}
\begin{loesungkarte}{gruen}{1.59}{8051 nach Fermat}
\derweg{$a$ von $\lceil\sqrt n\rceil$ aufwärts, bis $a^{2} - n$ ein Quadrat ist.}
\schritt{(a)}
$\sqrt{8051} = 89{,}7$; $a = 90$: $8100 - 8051 = 49 = 7^{2}$. Also $8051 = (90 - 7)(90 + 7) = 83\cdot97$. Ein Schritt.
\schritt{(b)}
Fermat findet die Zerlegung $n = pq$ bei $a = \frac{p + q}2$. Für $8057 = 7\cdot1151$: $a = 579$; von $\lceil\sqrt{8057}\rceil = 90$ bis $579$ sind es $490$ Schritte. Je weiter die Faktoren auseinander liegen, desto länger dauert es --- im schlimmsten Fall ($p$ klein) etwa $n/2$ Schritte, schlechter als Probedivision.
\schritt{(c)}
$p$ und $q$ dürfen nicht nahe beieinander liegen: Ist $|p - q| < n^{1/4}$ (etwa), findet Fermat sie sofort. Die Standards verlangen, dass sich $p$ und $q$ in den oberen Bits unterscheiden. (Umgekehrt dürfen sie auch nicht zu verschieden sein, sonst hilft die elliptische-Kurven-Methode, die kleine Faktoren schnell findet.)
\begin{mitnehmen}\kastentext
Fermat: $n = a^{2} - b^{2}$ --- schnell für nahe Faktoren, nutzlos sonst. Die Idee \enquote{finde $x^{2}\equiv y^{2}\pmod n$ mit $x\not\equiv\pm y$} ist aber die Grundlage der besten Verfahren (quadratisches Sieb, Zahlkörpersieb), die mit Kongruenzen statt Gleichungen arbeiten.
\end{mitnehmen}
\end{loesungkarte}

%%%%%%%%%% 1.60
\stupserzwei{gruen}{$x\colon2, 5, 26, 677$; $y\colon2, 26, 7474, 871$. $\operatorname{ggT}(|677 - 871|, 8051) = \operatorname{ggT}(194, 8051) = 97$.}{22pt}
\begin{loesungkarte}{gruen}{1.60}{8051 nach Pollard}
\derweg{Hase und Igel auf der Folge $x\mapsto x^{2} + 1$; der ggT der Differenz mit $n$ wird irgendwann größer als $1$.}
\schritt{(a)}
$f(x) = x^{2} + 1\bmod8051$, $x_0 = y_0 = 2$.
\begin{center}
{\small\renewcommand{\arraystretch}{1.3}
\begin{tabular}{@{}c|ccc@{}}
Schritt & $x_i = f(x_{i-1})$ & $y_i = f(f(y_{i-1}))$ & $\operatorname{ggT}(|x_i - y_i|, 8051)$\\\hline
1 & $5$ & $26$ & $1$\\
2 & $26$ & $7474$ & $1$\\
3 & $677$ & $871$ & $\operatorname{ggT}(194, 8051) = 97$\\
\end{tabular}}
\end{center}
$8051 = 97\cdot83$. Drei Schritte.
\schritt{(b)}
Modulo $97$: $x\colon2, 5, 26, 677\equiv95$; $y\colon2, 26, 7474\equiv5, 871\equiv95$. Im dritten Schritt ist $x_3\equiv y_3\pmod{97}$, also $97\mid x_3 - y_3$ --- aber modulo $83$ sind sie verschieden ($677\equiv13$, $871\equiv41$), also $8051\nmid x_3 - y_3$, und der ggT ist genau $97$. Die Folge modulo $97$ läuft in einen Zyklus, lange bevor sie es modulo $8051$ tut: Nach dem Geburtstagsparadoxon (Aufgabe 4.8) nach etwa $\sqrt{97}\approx10$ Schritten.
\schritt{(c)}
Die Folge modulo $p$ hat einen Vorlauf und dann einen Zyklus --- gezeichnet sieht das aus wie der Buchstabe $\rho$. Floyds Trick (Hase $y$ doppelt so schnell wie Igel $x$) findet den Zyklus, ohne die Folge zu speichern.
\begin{mitnehmen}\kastentext
Pollard $\rho$ findet einen Faktor $p$ in etwa $\sqrt p$ Schritten --- gut für Faktoren bis etwa $10^{12}$; für RSA-Zahlen ($p\approx10^{300}$) hoffnungslos. Die Idee, modulo eines unbekannten $p$ zu rechnen und den ggT den Rest tun zu lassen, ist eine der schönsten der Algorithmik.
\end{mitnehmen}
\end{loesungkarte}

%%%%%%%%%% 1.61
\stupserzwei{akzent}{$(5\cdot2^{7})^{4} = 5^{4}\cdot2^{28}\equiv1$ und $5^{4}\equiv-2^{4}$, also $-2^{4}\cdot2^{28} = -2^{32}\equiv1$.}{22pt}
\begin{loesungkarte}{akzent}{1.61}{Fermats Vermutung über $2^{2^{n}} + 1$}
\derweg{Zwei Darstellungen von $641$, zu Kongruenzen gemacht und kombiniert; dann die Faktorisierung von $2^{ab} + 1$; dann die Moral.}
\schritt{(a)}
$641 = 5\cdot2^{7} + 1$, also $5\cdot2^{7}\equiv-1\pmod{641}$, also $5^{4}\cdot2^{28}\equiv1$. $641 = 5^{4} + 2^{4}$, also $5^{4}\equiv-2^{4}$. Einsetzen: $-2^{4}\cdot2^{28} = -2^{32}\equiv1$, also $2^{32}\equiv-1$, also $641\mid2^{32} + 1 = F_5 = 4\,294\,967\,297$. $\blacksquare$ (Der Kofaktor ist $6\,700\,417$, prim.)
\schritt{(b)}
Für ungerades $b$: $x^{b} + 1 = (x + 1)(x^{b-1} - x^{b-2} + \dots + 1)$. Mit $x = 2^{a}$: $2^{ab} + 1$ ist durch $2^{a} + 1$ teilbar. Hat $m$ einen ungeraden Teiler $b > 1$, $m = ab$, so ist $2^{a} + 1$ ein echter Teiler von $2^{m} + 1$. Also: $2^{m} + 1$ prim $\Rightarrow$ $m$ hat keinen ungeraden Teiler $\Rightarrow$ $m = 2^{n}$. $\blacksquare$
\schritt{(c)}
Fünf Beispiele sind kein Beweis; und ein Muster, das bei kleinen Zahlen gilt, kann bei der nächsten brechen. Eulers Argument zeigt auch, \emph{wie} man Gegenbeispiele findet: Teiler von $F_n$ haben die Form $k\cdot2^{n+1} + 1$ (Analogon zu Aufgabe 1.45c), für $F_5$ also $64k + 1$ --- $641 = 10\cdot64 + 1$ ist der zweite Kandidat nach $193$. Heute weiß man: $F_5$ bis $F_{32}$ sind alle zusammengesetzt, und es ist offen, ob es überhaupt eine sechste Fermat-Primzahl gibt. Die Vermutung des besten Zahlentheoretikers seiner Zeit war falsch --- und das Nachrechnen hat es entschieden.
\begin{fehler}\kastentext
\enquote{Die ersten fünf sind prim, also alle.} --- Der häufigste Fehler der experimentellen Mathematik. Gegenbeispiele können sehr spät kommen: Die Vermutung $\pi(x) < \operatorname{li}(x)$ gilt für alle geprüften $x$ und ist trotzdem falsch (Littlewood 1914; erstes Gegenbeispiel unter $10^{317}$, Skewes).
\end{fehler}
\begin{mitnehmen}\kastentext
Eine Zahlentheorie-Vermutung prüft man an der Struktur, nicht an der Länge der Beispielliste: Welche Form müssen Teiler haben? Gibt es eine Faktorisierung? Eulers Teiler $641$ kam aus genau dieser Frage.
\end{mitnehmen}
\end{loesungkarte}
